\documentclass[10pt]{article}
\usepackage[letterpaper,margin=1in]{geometry}
\usepackage[T1]{fontenc}
\usepackage{lmodern,microtype,booktabs}
\usepackage{amsmath,amssymb,amsthm,mathtools,mathrsfs,hyperref}
\hypersetup{hidelinks,pdftitle={Bounds for the Bilu--Linial signing constant}}
\numberwithin{equation}{section}
\allowdisplaybreaks[1]
\setlength{\emergencystretch}{2em}
\newtheorem{theorem}{Theorem}
\newtheorem{lemma}{Lemma}[section]
\newtheorem{proposition}[lemma]{Proposition}
\newtheorem{corollary}[lemma]{Corollary}
\newcommand{\E}{\mathbb E}
\newcommand{\Pp}{\mathbb P}
\newcommand{\tr}{\operatorname{tr}}
\newcommand{\diag}{\operatorname{diag}}
\newcommand{\Var}{\operatorname{Var}}
\newcommand{\Cov}{\operatorname{Cov}}
\newcommand{\1}{\mathbf1}
\newcommand{\R}{\mathbb R}
\newcommand{\cP}{\mathcal P}
\newcommand{\cK}{\mathcal K}
\newcommand{\cD}{\mathcal D}
\newcommand{\norm}[1]{\left\lVert#1\right\rVert}
\begin{document}
\begin{center}
{\Large Bounds for the Bilu--Linial signing constant}
\end{center}
For a finite simple graph $G$, let $A_\sigma(G)$ be its signed adjacency
matrix and define
\[
 \gamma_d=\sup_{\Delta(G)\le d}\;
       \min_{\sigma:E(G)\to\{-1,1\}}\|A_\sigma(G)\|,
\]
where $\|\cdot\|$ is the Euclidean operator norm.
\begin{theorem}\label{thm:bounds}
There is an absolute constant $C$ such that, for all sufficiently large
integers $q$,
\[
 2\sqrt q+2q^{-5/2}
 -\left(\frac{32P_*}{\sqrt3}+o(1)\right)q^{-3}
 \le\gamma_{q+1}\le\sqrt{4q+Cq^{-2/17}}.
\]
\end{theorem}
Here $P_*\approx0.7632$ is the Sherrington--Kirkpatrick ground-state
constant, and $o(1)$ tends to zero as $q\to\infty$.

\section{Upper bound}

\subsection{Sources, deletion, and the maximum principle}
Put $q=d-1$, $R^2=4q+\Delta$, $a=R^{-1}$, and define
\[
 \frac{q\eta_0^2}{1-\eta_0}=\Delta,\qquad
 \tau_*=(1-\eta_0)/q,\qquad
 s=\frac{1+q\tau_*}{1-\tau_*},\qquad r=1+\eta_0/2.
\]
In the application $\Delta=4/p$ and
$p=\lfloor c_0d^{2/17}\rfloor$, where $0<c_0\le1$ is an absolute
constant chosen below.
For a graph $H$ of maximum degree $d$ and $0<y_i\le s$, set
\[
 c_{ij}(1+c_{ij})=a^2y_iy_j,\quad
 \tau_{ij}=\frac{c_{ij}}{1+c_{ij}},\quad
 Z_i(y)=\frac{1+\sum_{j\sim i}c_{ij}}{y_i},\quad D_i=y_iZ_i.
\]
Then $\tau_{ij}\le\tau_*$, $s\le2+C/d$,
$da^2s^2\le1+C/d$, and $c_{ij}\le C/d$.
Use independent sources $y^+,y^-$ and precisions
$P^\pm=\operatorname{diag}Z(y^\pm)\pm aA_\sigma$.
The sign law has weight
$W(\sigma)=\mathbf1_{\{P^\pm\succ0\}}(\det P^+\det P^-)^p$.
Write $G^\pm=(P^\pm)^{-1}$ and $h_i^\pm=G^\pm_{ii}/y_i^\pm$.
Congruence by $\sqrt{\operatorname{diag}y^\pm}$ multiplies all
weights by the same source-dependent factor and extends the law
to zero sources. A zero source is an isolated identity row in the
normalized matrix, with normalized inverse diagonal one and physical
row zero. All ensuing formulas at zero sources mean these limits.

\begin{lemma}[Deletion and source differentiation]\label{lem:fw-deletion}
Writing $D_y=\operatorname{diag}y$, one has
\[
 DZ=-D_y^{-1}\mathcal B D_y^{-1},\quad
 \mathcal B_{ii}=1+\sum_j\frac{\tau_{ij}^2}{1-\tau_{ij}^2},
 \quad \mathcal B_{ij}=-\frac{\tau_{ij}}{1-\tau_{ij}^2}\ (i\sim j).
\]
The matrix $\mathcal B$ is positive definite, and
$K=\mathcal B^{-1}$ is entrywise nonnegative, with
$K_{ii}\ge1$ and $K_{ij}\ge\tau_{ij}$ on edges.
For every principal deletion $J=H-v$, there is a deterministic
$\widehat y\le y|_J$ satisfying $Z^J(\widehat y)=Z^H(y)|_J$.
\end{lemma}
\begin{proof}
Differentiate $c(1+c)=a^2y_iy_j$. The displayed Jacobian follows,
and its diagonal-dominance margin is
\[
 1-\sum_j\frac{\tau_{ij}}{1+\tau_{ij}}
 \ge1-\frac{d\tau_*}{1+\tau_*}
 =\frac{\eta_0}{1+\tau_*}>0.
\]
Thus $\mathcal B$ is a positive definite matrix with nonpositive
off-diagonal entries and nonnegative inverse.
To increase $Z(y)$ by $u\ge0$, lift the segment $Z(y)+tu$:
$y'(t)=-D_yKD_yu\le0$.
The lift remains bounded away from zero since
$Z_i(y(t))\ge1/y_i(t)$ implies
$y_i(t)\ge[Z_i(y)+u_i]^{-1}$; compact continuation reaches $t=1$.
The symmetric negative-definite Jacobian makes $Z$ injective by
integration along segments. Deletion increases each surviving
precision above $Z^J(y|_J)$ by $\mathbf1_{i\sim v}c_{iv}/y_i$,
so the lift proves the assertion.

For completeness, $K_{ij}$ is the sum of products of $\tau_e$
over nonbacktracking walks from $i$ to $j$, including the empty
walk. This series converges: the directed walk operator applied
to $f_{i\to j}=(1+\tau_{ij})^{-1}$ is strictly less than $f$ by
the same dominance inequality. If $S_{iw}$ is its walk sum and
$M_{i\to j,w}$ the part starting with $i\to j$, then
$S_{iw}=\mathbf1_{i=w}+\sum_jM_{i\to j,w}$ and
$M_{i\to j,w}=\tau_{ij}(S_{jw}-M_{j\to i,w})$.
Elimination gives $\mathcal B S_{\cdot w}=e_w$, proving the stated
lower bounds. Treat the two source branches separately.
\end{proof}

For a deleted core $J=H-v$ and $N=N_H(v)$, use its inherited
precisions, equivalently the smaller sources above. Put
\[
 A=\frac{a^2}{Z_v^+}G_J^+[N,N],\quad
 B=\frac{a^2}{Z_v^-}G_J^-[N,N],\quad
 \alpha=(1-\xi^TA\xi)_+,\quad\beta=(1-\xi^TB\xi)_+,
 \quad\Phi=\alpha^p\beta^p,
\]
where $\xi$ is a fresh uniform sign vector before insertion.
Let $f=\mathbb E_\xi\Phi$ and $F_H=\mathbb E_J f$.
The actual full-law core marginal is $f/F_H$ times its own law.
A bad core has no good extension. On a good extension,
\[
 Z_v^+G^+_{vv}=\alpha^{-1},\quad Z_v^-G^-_{vv}=\beta^{-1},\quad
 G^+_{vi}=-a^{-1}(A\xi)_i/\alpha,\quad
 G^-_{vi}=a^{-1}(B\xi)_i/\beta.                 \tag{F1}\label{eq:fw-schur}
\]

\begin{lemma}[Uniform insertion floor]\label{lem:fw-floor}
For sufficiently large $d$, if the proper-core means obey
$\mathbb E G_{J,ii}^\pm\le1.01y_i^\pm$, then $F_H\ge40^{-p}$.
Consequently the full law exists and
$\mathbb E(h_i^\pm)^k\le40^p$ for $1\le k\le p$.
\end{lemma}
\begin{proof}
Set $\kappa=a^2s^2$, $x_\pm=y_v^\pm/s$, $u_i^\pm=y_i^\pm/s$,
$C_\pm=\kappa x_\pm\sum_Nu_i^\pm$, and $t_\pm=C_\pm/D_v^\pm$.
For large $d$, $d\kappa\le1.01$, $c_{ij}\le.01$, and
\[
 C_\pm\le1.01,\quad D_v^\pm\ge1+C_\pm/1.01,
 \quad t_\pm\le.505,\quad
 \mathbb E(q_A+q_B)\le1.01(t_++t_-),
\]
where $q_A=\xi^TA\xi$, $q_B=\xi^TB\xi$.
The parallel-sum formula gives pointwise
\[
 q_Aq_B\ge h(q_A+q_B),\qquad
 h=\min_\pm\frac{a^2y_v^\pm}{D_v^\pm}
                  \sum_N\frac1{Z_i^++Z_i^-}.
\]
Indeed $P_J^++P_J^-$ is diagonal, and for positive matrices $P,Q$
the variational formula bounds
$\xi^T(P+Q)^{-1}\xi$ by $ab/(a+b)$, where
$a=\xi^TP^{-1}\xi$, $b=\xi^TQ^{-1}\xi$; rescaling gives the claim.
If $t_++t_-\le.95$, then
$\mathbb E\alpha\beta\ge1-1.01(.95)>1/40$.
Otherwise $C_\pm>.79$, $x_\pm>.78$,
$x_\pm/D_v^\pm>.38$, and
$\sum_N(u_i^++u_i^-)-|N|>.57/\kappa$.
Using $D_i^\pm\le2.01$ and
$uv/(u+v)\ge(u+v-1)/2$ on $[0,1]^2$ gives
$h>.38\cdot.57/(2\cdot2.01)>1/20$.
The product constraint therefore implies
$\alpha\beta\ge[1-.95(q_A+q_B)]_+$:
if one energy exceeds one, their sum is at least $20/19$;
otherwise expand their product.
Hence $\mathbb E\alpha\beta\ge1-.95(1.01)^2>1/40$ in either case.
Jensen proves the floor. Finally $h_v=(D_v\alpha)^{-1}$,
$D_v\ge1$, so multiplying the insertion weight by $h_v^k$
leaves nonnegative residual powers and a numerator at most one.
Deleting each root proves the moment bound.
\end{proof}

\begin{lemma}[Source maximum and moments]\label{lem:fw-maximum}
Suppose the law exists throughout a capped source cube and all
normalized first means are at most $B_0\ge1$. Then, for
$1\le k\le p-2$,
\[
 \sup\mathbb E(h_i^\pm)^k
       \le\left(\frac{pB_0-k}{p-k}\right)^k.        \tag{F2}\label{eq:fw-moments}
\]
At a maximum $\mathbb E h_v^+=r$ of a normalized first mean,
\[
 p\operatorname{Cov}(G^+_{vv},G^+_{ii})
       -\mathbb E(G^+_{vi})^2
       \le-r y_v^+y_i^+K_{vi}.                    \tag{F3}\label{eq:fw-covariance}
\]
In particular $(p-1)\operatorname{Var}(h_v^+)\le r(r-1)$.
\end{lemma}
\begin{proof}
In physical precision coordinates, write
$C_i=\partial_{Z_i}\mathbb E G_{vv}^k
=p\operatorname{Cov}(G_{vv}^k,G_{ii})
-k\mathbb E(G_{vv}^{k-1}G_{vi}^2)$.
At a maximum $m_k$ of $\mathbb E h_v^k$, decreasing each positive
source is feasible, so
$(DZ)^TC\ge k m_k y_v^{k-1}e_v$.
Multiply by $-D_yKD_y$ to obtain
$C_i\le-km_k y_v^k y_iK_{vi}$.
At $i=v$ this yields
$(p-k)\mathbb E h_v^{k+1}\le m_k(pB_0-k)$.
Moment log-convexity proves (F2); $k=1$ gives (F3) and its variance
consequence, since $K_{vv}\ge1$.
The finite sign sums and the floor justify maxima and derivatives:
$p-k\ge2$ makes inverse-weighted determinant powers continuously
differentiable across a support boundary. A zero-source root has
moment one, and other zero coordinates can be removed first.
\end{proof}


The induction uses a cube expanding from zero. Proper graphs satisfy
the mean cap; deletion and the floor make the new law positive on
its whole source cube before any cap is asserted there. A first
contact supplies (F2)--(F3), with bounded $L^k$ bases through $k\le p/2$.
Ruling out that contact therefore completes the graph-order induction.
At the final constant sources $y_i=s$,
\[
 Z_i(s\mathbf1)=\frac{1+(\deg(i)-1)\tau_*}{1+q\tau_*}\le1.
\]
Positivity of both $P^\pm$ then implies $I\pm A_\sigma/R\succ0$
for a supported signing.

\subsection{Endpoint comparison and concentration}

All expectations below use the paired law just defined. Constants $C,c>0$
are absolute and may change. We take
\[
 p=\lfloor c_0d^{2/17}\rfloor,\qquad
 k_*=\left\lceil16p/\log d\right\rceil,
\]
and work under the first-mean cap $\mathbb Eh_i^\pm\le r=1+\varepsilon$.
The source maximum bounds the $L^j$ norms of physical diagonals and of
$Z_iG_{ii}=D_ih_i$ by an absolute constant for $j\le p/2$.
Thus all moments of order $8k_*+12$ used below are available for large $d$.
For a root $v$, write $N=N(v)$, $K=H-v$ and
\[
 A=\frac{a^2}{Z_v^+}G_K^+[N,N],\quad
 B=\frac{a^2}{Z_v^-}G_K^-[N,N],\quad
 \alpha=(1-\xi^TA\xi)_+,\quad\beta=(1-\xi^TB\xi)_+,\quad
 \Phi=\alpha^p\beta^p.
\]
The own core law is denoted by $\nu_K$; its insertion factor is
$F_H=\mathbb E_{\nu_K,\xi}\Phi\ge40^{-p}$. The actual core marginal is
$f/F_H$ times $\nu_K$, where $f=\mathbb E_\xi\Phi$.
Set $q_M(x)=x^TMx$, $b_i=A_{ii}+B_{ii}$ and $\mu_0=\operatorname{tr}(A+B)$.

\begin{lemma}[Endpoint calculus]\label{cu:endpoint}
For an incident coordinate $\xi_i$, put
$\mathcal D_j\psi=W^{-1}\partial_i^j(W\psi)$, with deterministic
precisions held fixed. Then
\[
 \mathbb E[\xi_iG^+_{vi}]
 =\mathbb E[\mathcal D_1G^+_{vi}-\tfrac13\mathcal D_3G^+_{vi}]
       +O(p^5a^5).
 \tag{E1}\label{cu:E1}
\]
The error is not divided by an insertion factor.
\end{lemma}
\begin{proof}
On a good endpoint of the two-point edge fiber put
$D=G^+_{vv}+G^+_{ii}+G^-_{vv}+G^-_{ii}$ and $M=(64pa)^{-1}$.
Call the fiber regular if both endpoints are good and have $D\le M$.
A good endpoint with $D\le M/2$ lies on such a fiber: the rank-two
perturbation has relative norm at most $2aD$, so throughout the segment
$G^\pm(s)\preceq2G^\pm(\xi_i)$ and $W(s)\le e^{1/16}W(\xi_i)$.
At every good point,
\[
 \partial_iG^\pm_{jl}=\mp a(G^\pm_{jv}G^\pm_{il}+G^\pm_{ji}G^\pm_{vl}),
 \qquad \partial_i\log W=2pa(G^+_{vi}-G^-_{vi}).
 \tag{E3}\label{cu:E3}
\]
Consequently $|\mathcal D_jG^+_{vi}|\le C_jp^ja^jD^{j+1}$.
Taylor's endpoint identity
$\mathbb E_\xi\xi f(\xi)=\mathbb E_\xi(f'(\xi)-f'''(\xi)/3)
+O(\|f^{(5)}\|_\infty)$ is exact through degree four. Applied on regular
fibers, averaged before dividing by the full normalizer, it costs
$Cp^5a^5\mathbb ED^6$. On every omitted good endpoint $D>M/2$;
for $j=0,1,3$, its removed contribution is at most
$C_jp^ja^j(2/M)^{5-j}\mathbb ED^6$. This proves \eqref{cu:E1}.
\end{proof}

We also need a high-order comparison whose retained terms are evaluated
at sign endpoints. Define $c_j$ by
\[
 e^{t^2/2}/\cosh t=1+\sum_{j\ge2}c_jt^{2j},
 \qquad |c_j|\le C^j,\quad c_2=1/12.
\]
For a multi-index $\mathbf j$ with entries zero or at least two, put
$|\mathbf j|_g=\sum_{j_i>0}(j_i-1)$ and $c_{\mathbf j}=\prod c_{j_i}$.
If $\mathsf G,\mathsf R$ denote Gaussian and sign expectation, then
\[
 \mathsf Gf=\sum_{|\mathbf j|_g\le k}c_{\mathbf j}
       \mathsf R\partial^{2\mathbf j}f+\mathcal E_k(f),\qquad
 |\mathcal E_k(f)|\le C^{k+1}
       \sum_{|\mathbf j|_g=k+1}\|\partial^{2\mathbf j}f\|_\infty.
 \tag{E4}\label{cu:E4}
\]
The sum in the remainder uses distinct active coordinates, and every
derivative has order at most $4k+4$. To prove this, in one coordinate
Taylor-expand through degree $2L+3$. Matching coefficients of
$e^{t^2/2}/\cosh t$ leaves a remainder at most
$C^{L+1}\|\partial^{2L+4}f\|_\infty$. Process the coordinates in order,
spending grade $j-1$ on a selected derivative $2j$. A first remainder
has exact total grade $k+1$; the number of grade compositions is at most
$2^k$. This proves \eqref{cu:E4}, also for bounded weak highest derivatives
by smoothing. Clipped powers below have the required regularity.

Here is the derivative accounting used with \eqref{cu:E4}. On common
residual support, the normalized atoms $(Ax)_i/\sqrt{b_i}$ and
$(Bx)_i/\sqrt{b_i}$ and their normalized derivatives are bounded by one.
After $n$ derivatives the coefficient mass grows by at most $4p+n+1$.
Thus, for $n\le4k_*+4$,
\[
 \begin{split}
 |\partial^\nu[(\operatorname{tr}A-q_A)\alpha^{p-1}\beta^p]|
   &\le(Cp)^n(2+\mu_0)\prod_i b_i^{\nu_i/2},\\
 |\partial^\nu F|&\le p(Cp)^n\mu_0\prod_i b_i^{\nu_i/2},\quad
 F=(p-1)q_{A^2}\alpha^{p-2}\beta^p
                 +pq_{AB}\alpha^{p-1}\beta^{p-1}.
 \end{split} \tag{E5}\label{cu:E5}
\]
The quadratic prefactors, divided by $2+\mu_0$ or $\mu_0$, and their
first two normalized derivatives are bounded; higher ones vanish.
These facts prove the bounds by induction. A grade-$k+1$ remainder
with $l$ active coordinates has total half-order $k+1+l$.
Since $b_i\le Ca^2(G_{K,ii}^++G_{K,ii}^-)$, core moments and H\"older
bound its average by $C^{k+l+1}d^{-k-1-l}$; the $d^l$ choices cancel
the last power. The averaged remainders in \eqref{cu:E5} are therefore
$(Cp)^{4k+4}d^{-k-1}$ and $p(Cp)^{4k+4}d^{-k-1}$.

At actual sign endpoints, division by $\Phi$ and \eqref{cu:E3} instead
give coefficient mass $(Cp)^n$, a factor $a^n$, and inverse-entry
degree $n$. For the first expression of \eqref{cu:E5}, use exactly
\[
 \frac{(\operatorname{tr}A-q_A)\alpha^{p-1}\beta^p}{\Phi}
   =1+(\operatorname{tr}A-1)t,\qquad t=Z_v^+G^+_{vv},
\]
and keep $t$ grouped: $\partial_i t=-2atG^+_{vi}$.
Also $\operatorname{tr}A\le Ca^2\sum_iG^+_{ii}$ by Schur order.
Diagonal moments then bound a retained grade-$j$ term by
$(Cp)^{4j}d^{-j}$. For $F/(a^2\Phi)$ there is one free row index,
inverse degree $n+2$, and coefficient mass $p(Cp)^n$; its grade-$j$
term has $L^2$ norm at most $p(Cp)^{4j}d^{1-j}$.
These statements follow directly from
$|G_{ij}|\le\sqrt{G_{ii}G_{jj}}$ and use moments at most $8j+4$.
All derivatives vanish at excluded endpoints, so passing to the actual
law is exact. Finally,
\[
 p40^p(Cp)^{4k_*+4}d^{-k_*}\le e^{-cp},
 \tag{E6}\label{cu:E6}
\]
because $p\le d^{1/7}$ makes its logarithm at most
$[\log40-16(1-4/7)+o(1)]p<0$.

\begin{lemma}[Bias and reference concentration]\label{cu:concentration}
Put $N_R=\mathsf R[(\operatorname{tr}A-q_A)\alpha^{p-1}\beta^p]$
and define $N_G$ with Gaussian expectation. Uniformly throughout the
capped source family,
\[
 \left|\mathbb E_{\nu_K}(N_R-N_G)/F_H\right|\le Cp^4/d,
 \qquad N_G\ge0. \tag{C1}\label{cu:C1}
\]
Moreover
\[
 \begin{split}
 &\sup_{i,\pm}\mathbb E(h_i^\pm-1)^2
  +\sup_{i,\pm}(1-\mathbb Eh_i^\pm)_+
  +\sup_{v,\pm}a^2\sum_{i\sim v}\mathbb E(G^\pm_{vi})^2\\
 &\hspace{12mm}
  +\sup_v\mathbb E_{\rm actual\ core}\operatorname{tr}(A^2+B^2)
       \le\delta:=C\{\varepsilon+p^4/d+(p/d)^{1/3}\}.
 \end{split} \tag{C2}\label{cu:C2}
\]
\end{lemma}
\begin{proof}
Apply \eqref{cu:E4}--\eqref{cu:E6} to the first expression in
\eqref{cu:E5}: grade one costs $Cp^4/d$, higher grades sum geometrically
to $Cp^8/d^2$, and the only floor loss is exponentially small.
For Gaussian nonnegativity, an even log-concave density has covariance
at most $I$ relative to the standard Gaussian: its convolution with
a translated Gaussian is even log-concave by Pr\'ekopa's theorem~\cite{Prekopa},
so its Hessian at zero is nonpositive. Apply this to
$\alpha^{p-1}\beta^p$ and take its covariance trace against $A$.

Write $\rho=\sup(1-\mathbb Eh_i)_+$ and
$b=\varepsilon+Cp^4/d+d^{-1}$. The exact Schur identity is
\[
 \mathbb E_{\nu_K}N_R/F_H
 =1-D_v\mathbb Eh_v
   +a^2\sum_{i\sim v}\mathbb E(G_{vv}G_{ii}-G_{vi}^2).
 \tag{C3}\label{cu:C3}
\]
Since $D_v=1+L_v-C_v$, $L_v=a^2y_v\sum_i y_i\le1+C/d$,
$C_v=\sum_i c_{vi}^2=O(d^{-1})$, and $\mathbb Eh_i^2\le1+C\varepsilon$,
\eqref{cu:C1} gives the row bound
$a^2\mathbb E\sum_iG_{vi}^2\le C(\rho+b)$.

The source second-moment bound already gives
$\mathbb E(h_i-1)^2\le C\varepsilon+2\rho$.
For the core kernel put $T_A=\operatorname{tr}A^2$,
$\mu=\operatorname{tr}A$ and $w=T_A/(1+\mu)\le\mu$.
Endpoint expansion of the core-constant numerator $w\Phi$ gives
$|\mathbb E_{\nu_K}w(f_G-f_R)/F_H|\le Cp^4/d$.
Whiten the Gaussian potential by $I+2(p-1)A$; its remaining factor
is even log-concave because $-\log(1-q_A)-q_A$ is convex.
The covariance bound therefore gives
\[
 N_G\ge f_G\frac{2(p-1)T_A}{1+2(p-1)\|A\|}\ge f_Gw.
\]
Dropping the negative row term in \eqref{cu:C3} yields
$\mathbb Ew\le C(\rho+b)$.
Also $\mu\le U=(a^2/Z_v)\sum_iG_{ii}$ and
$U_0=(a^2/Z_v)\sum_i y_i=L_v/D_v\le1/2+C/d$.
For $k=\lfloor p/4\rfloor$, Minkowski and the source moment caps give
$\|U\|_k\le U_0(pr-k)/(p-k)\le3/5$ for sufficiently large $d$.
Consequently
\[
 \mathbb E[T_A1_{\{\mu>1\}}]
 \le\mathbb E[U^2 1_{\{U>1\}}]\le\mathbb EU^k\le e^{-cp},
 \qquad T_A1_{\{\mu\le1\}}\le2w.
 \tag{C3a}\label{cu:core-tail}
\]
This uses only the capped moments and is uniform at zero sources.
Applying the same argument to both branches proves the preclosure bound
\[
 \mathbb E\Theta\le C(\rho+b),\qquad
 \Theta=\operatorname{tr}(A^2+B^2).
 \tag{C3b}\label{cu:preclosure-core}
\]

Let $X_z=(P+zI)^{-1}$ and
$t=\sup_{i,\pm}(1-\mathbb E X_{z,ii}^\pm/y_i^\pm)_+$.
Then $\rho\le t$ and
$\mathbb E(G_{ii}^\pm-X_{z,ii}^\pm)\le y_i^\pm(\varepsilon+t)$.
For a root $v$, write $C_\nu=(P^\nu_{-v})^{-1}$ and
$C_{z,\nu}=(P^\nu_{-v}+zI)^{-1}$; restrictions to $N=N(v)$
will be indicated when needed. Put
\[
 M=\frac{a^2}{Z_v^+}C_{z,+}[N,N]\preceq A.
\]

We need a lower bound on the covariance of a Gaussian under the
fresh-star weight $\Phi=\alpha^p\beta^p$. For this normalized Gaussian
law, let $\Sigma$ be its covariance and
$V=-p\log\alpha-p\log\beta$ on its convex support.
Integration by parts gives
$\mathbb E[x(x+\nabla V)^T]=I$ and
$\mathbb E[(x+\nabla V)(x+\nabla V)^T]=I+\mathbb E\nabla^2V$.
Matrix Cauchy--Schwarz therefore gives
$\Sigma\succeq(I+\mathbb E\nabla^2V)^{-1}$.
The boundary terms vanish because $p$ is large; the formulas can
also be obtained by smoothing the clipped density.
Since $\nabla^2V\succeq0$ and $I-(I+Q)^{-1}\preceq Q$ for $Q\succeq0$,
\[
 \begin{split}
 \mathsf G[(\operatorname{tr}M-q_M)\Phi]
 \le{}&2p\,\mathsf G[(\operatorname{tr}(MA)\alpha^{-1}
                         +\operatorname{tr}(MB)\beta^{-1})\Phi]\\
 &+4p\,\mathsf G[(q_{AMA}\alpha^{-2}
                         +q_{BMB}\beta^{-2})\Phi].
 \end{split}\tag{C4}\label{cu:C4}
\]
This follows from
$\nabla^2(-\log\alpha)=2A/\alpha+4Axx^TA/\alpha^2$ and its
$B$ analogue. All four terms on the right are nonnegative.

We transfer \eqref{cu:C4} to signs using \eqref{cu:E4}.
All four right-hand prefactors retain the core factor $\Theta$:
\[
 \operatorname{tr}(MA),\operatorname{tr}(MB)\le\Theta,
 \qquad AMA\preceq\Theta A,\quad BMB\preceq\Theta B.
\]
For either quadratic kernel $L$ and the corresponding $Q=A$ or $B$,
residual support gives
\[
 q_L\le\Theta,\quad |\partial_iq_L|\le2\Theta\sqrt{Q_{ii}},
 \quad |\partial_i\partial_jq_L|\le2\Theta\sqrt{Q_{ii}Q_{jj}}.
\]
More explicitly, each unmultiplied right-hand function $F$ satisfies
at a sign endpoint
\[
 |\partial_i^4F|/\Phi
 \le Cp^4a^4\Theta(1+D_i)^C(1+t_++t_-)^2,
 \quad t_\nu=Z_v^\nu G^\nu_{vv},
\]
where $D_i$ is one plus the four physical inverse diagonals at $v,i$.
The product rule and the identities
$\alpha'/\alpha=2aG^+_{vi}$,
$\alpha''/\alpha=-2a^2(G^+_{vv}G^+_{ii}-(G^+_{vi})^2)$,
and their minus counterparts prove this bound.
By \eqref{cu:preclosure-core}, $\rho\le t$ and high-moment
interpolation, the four right-hand grade-one corrections, including
their factors $p$, total $Cp^5(t+b)/d$.

For the left function $F_0=(\operatorname{tr}M-q_M)\Phi$, put
$R_i=|G^+_{vi}|+|G^-_{vi}|$.
A first residual derivative supplies a root row entry, a second
supplies two inverse entries, and all higher residual derivatives
vanish. Hence
\[
 \begin{aligned}
 |\Phi''|/\Phi&\le Ca^2(p^2R_i^2+pD_i^C),\\
 |\Phi'''|/\Phi&\le Ca^3(p^3R_i^3+p^2R_iD_i^C),\\
 |\Phi''''|/\Phi&\le Ca^4(p^4R_i^4+p^3R_i^2D_i^C+p^2D_i^C).
 \end{aligned}
\]
Since $q_M\le q_A\le1$, $M_{ii}\le Ca^2G^+_{ii}$,
$|\partial_iq_M|\le CaD_i$ and
$|\partial_i^2q_M|\le Ca^2D_i$. Thus
\[
 |\partial_i^4F_0|/\Phi
 \le Ca^4(1+\mu)D_i^C
 \{p^4R_i^4+p^3(R_i^2+R_i^3)+p^2(1+R_i+R_i^2)\}.
\]
For a uniform incident index $I$, padded to $d$ slots, the row bound
from \eqref{cu:C3} gives $\mathbb ER_I^2\le C(t+b)$.
Interpolation preserves this scale in every monomial with at least
two row factors; $t+b\ge b$ has a polynomial lower envelope.
The left grade-one correction is consequently
$C\{p^4(t+b)+p^2\}/d$.
The higher grades of the five functions retain their geometric
bounds, totaling $Cp^9/d^2$ after the outer factors $p$.
Their core prefactors have capped high moments, so the final
insertion-floor remainder is exponentially small as in \eqref{cu:E6}.
All source factors are dimensionless $t_\nu$ or physical inverse
entries, so these estimates extend to zero sources.

We next bound the four transferred terms. Write
\[
 W=\frac1d\sum_{\nu=\pm}\operatorname{tr}
                   [(C_\nu-C_{z,\nu})[N,N]]\ge0.
\]
The full/core shifted rank-one correction satisfies
\[
 \operatorname{tr}[X_z^\nu[N,N]-C_{z,\nu}[N,N]]
 =a^2(X_z^\nu)_{vv}\,\xi^TC_{z,\nu}[N,N]^2\xi
 \le\frac{Z_v^\nu}{z}(G^\nu_{vv}-(X_z^\nu)_{vv}).
 \tag{C5}\label{cu:C5}
\]
Indeed $C_\nu-C_{z,\nu}\succeq zC_{z,\nu}^2$, and compressing the
square dominates the square of the compression. The root Schur
identities give
$a^2\xi^T(C_\nu-C_{z,\nu})[N,N]\xi
=(X_z^\nu)_{vv}^{-1}-(G^\nu_{vv})^{-1}-z$.
Multiplication by $(X_z^\nu)_{vv}/z$ proves the inequality.
Since the unshifted full/core correction is nonnegative, it follows
that $\mathbb EW\le C(t+\varepsilon)$ whenever $dz\ge1$.
Also $\mathbb EW^2\le C$, since $W$ is bounded by a normalized sum
of full inverse diagonals. H\"older and the source moments give
\[
 \mathbb E[(h_v^++h_v^-)W]
 \le C(t+\varepsilon)^{1-O(1/p)}\le C(t+\varepsilon).
\]
The last step uses $t+\varepsilon\ge\varepsilon$, which is a negative
power of $d$, whereas $p$ is a positive power.

For either branch $\nu$, positivity gives
\[
 \begin{split}
 \operatorname{tr}(C_{z,+}[N,N]C_\nu[N,N])
 &\le\tfrac12\operatorname{tr}(C_{z,+}[N,N]^2+C_{z,\nu}[N,N]^2)
       +z^{-1}\operatorname{tr}(C_\nu-C_{z,\nu})[N,N]\\
 &\le CdW/z.
 \end{split}
\]
Consequently the expectation of the two trace terms in
\eqref{cu:C4}, after division by the full normalizer, is at most
$C(t+\varepsilon)/(dz)$: for example
$\operatorname{tr}(MA)/\alpha
=a^4G^+_{vv}\operatorname{tr}(C_{z,+}[N,N]C_+[N,N])/Z_v^+$.
Here $1/Z_v^+\le s$, and the preceding weighted bound on $W$ applies.
The quadratic terms are at most
\[
 a^2\|M\|\,\mathbb E\sum_i[(G^+_{vi})^2+(G^-_{vi})^2]
 \le C(t+b)/(dz),
\]
since $\|M\|\le Ca^2/z$ and the row bound was proved above.
Let $q_z=a^2\xi^TC_{z,+}[N,N]\xi$ and
$\mu_z=a^2\operatorname{tr}C_{z,+}[N,N]$.
Their normalized bias is therefore
\[
 \mathbb E(\mu_z-q_z)/Z_v^+
 \le Cp(t+b)/(dz)+C\{p^5(t+b)/d+p^2/d+p^9/d^2+e^{-cp}\}.
 \tag{C6}\label{cu:C6}
\]
In the root equation this is multiplied by $y_v^+Z_v^+=D_v^+\le C$.
The shifted core/full diagonal replacement, using \eqref{cu:C5},
costs only $C(t+\varepsilon)/(dz)$ after the same normalization.
To make the last scalar step explicit, put $m_i=\mathbb E(X_z)_{ii}$.
Jensen gives $m_v^{-1}\le Z_v+z-a^2\sum_i m_i+e_v$, where the
combined bias and core-replacement errors obey
$y_ve_v\le Cp(t+b)/(dz)+C\{p^5(t+b)/d+p^2/d+p^9/d^2+e^{-cp}\}$.
At a maximizing vertex $m_v=y_v(1-t)$ and $m_i\ge y_i(1-t)$.
Since $y_vZ_v=1+L_v-C_v$, subtraction yields
\[
 \frac{t^2}{1-t}\le(L_v-1)t-C_v+y_v(z+e_v).
\]
Using $L_v\le1+C/d$, $C_v\ge0$, and $y_v\le C$, we obtain
\[
 t^2\le C\{z+p(t+b)/(dz)+p^5(t+b)/d+p^2/d+p^9/d^2+d^{-1}+e^{-cp}\}.
\]
Take $u=(p/d)^{1/3}$ and $z=u^2$. Then
\[
 t^2\le C\{u^2+u(t+b)+p^5(t+b)/d+p^2/d+p^9/d^2+e^{-cp}\}.
\]
For $p\le d^{1/7}$ one has $p^5/d\le u$, $b\le Cu$,
$p^2/d\le u^2$ and $p^9/d^2\le u^2$.
Therefore $t\le Cu$. The row, second-moment and core bounds already
proved in terms of $\rho\le t$ give \eqref{cu:C2} directly.
\end{proof}

\begin{lemma}[Uniform incident-row gain]\label{lem:global-row-gain}
Under the capped source law,
\[
 \sup_{v,\nu}a^2\sum_{i\sim v}\mathbb E(G^\nu_{vi})^2
 \le\delta_{\rm row}:=C\delta/p.
 \tag{GR1}\label{eq:global-row-gain}
\]
This improves the incident-row envelope, while the core Frobenius
envelope remains $\delta$.
\end{lemma}
\begin{proof}
For PSD $A,B$, Royen's full Gaussian correlation monotonicity,
in the formulation of Milman~\cite[Section 3.3]{Milman}, gives
\[
 \mathsf G[q_{AB}\alpha^{p-1}\beta^{p-1}]\ge0.
\]
Indeed, if standard Gaussian vectors $X,Y_s$ have cross covariance
$sI$, then $H(s)=\mathbb E\alpha(X)^p\beta(Y_s)^p$ is nondecreasing
on $[0,1]$: both factors are nonnegative even quasi-concave functions.
Gaussian differentiation gives
$H'(1)=4p^2\mathsf G[q_{AB}\alpha^{p-1}\beta^{p-1}]$.
Bounded gradient products justify continuity at $s=1$; clipping is
sufficiently smooth because $p>2$.

Apply \eqref{cu:E4} to $F=q_{AB}\alpha^{p-1}\beta^{p-1}$.
Its prefactor and first two normalized derivatives are bounded by
$C\mu_0$, so on residual support
$|\partial^\nu F|\le(Cp)^{|\nu|}\mu_0\prod_i b_i^{\nu_i/2}$.
At sign endpoints,
\[
 F/(a^2\Phi)=-\sum_{j\sim v}G^+_{vj}G^-_{vj}.
\]
After $n$ derivatives its coefficient mass is $(Cp)^n$, with a
factor $a^n$, inverse degree $n+2$, and one free row index.
In grade $j$, with $l\le j$ active coordinates, $n=2(j+l)$ and
$a^nd^{l+1}=O(d^{1-j})$.
Capped diagonal moments bound the retained grade by
$(Cp)^{4j}d^{1-j}$ under the actual sign law.
The only insertion-floor loss is in the final remainder,
\[
 40^p(Cp)^{4k_*+4}d^{-k_*}\le e^{-cp},
\]
which holds throughout $p\le d^{1/7}$ since
$\log40-16(1-4/7)<0$. Hence
\[
 T_v:=\sum_{i\sim v}\mathbb E G^+_{vi}G^-_{vi}\le Cp^4.
\]

Write $m_v=\mathbb Eh_v^+$ and $S_v=\sum_i\mathbb E(G^+_{vi})^2$.
The diagonal resolvent equation and the endpoint identity give
\[
 (2p-1)a^2S_v=B_v+2pa^2T_v+O(p^3/d),\qquad
 B_v=1-D_v m_v+a^2\sum_i\mathbb E G^+_{vv}G^+_{ii}.
\]
The third derivative costs $Cp^3a^4d=Cp^3/d$;
the fifth-order remainder is $Cp^5a^6d=Cp^5/d^2$.
With $\rho=\sup(1-\mathbb Eh_i)_+$, the moment caps and
$D_v=1+L_v-C_v$ imply
\[
 B_v\le1-D_v(1-\rho)+L_v(1+C\varepsilon)
       \le C(\rho+\varepsilon+d^{-1}).
\]
Thus $a^2S_v\le C\{(\rho+\varepsilon+d^{-1})/p+p^4/d\}$.
Use $\rho\le\delta$ from (C2) and
$p^4/d\le(p/d)^{1/3}/p$ when $p^7\le d$.
The minus branch is identical. All derivatives involve physical
inverse entries, so the proof remains valid at zero sources.
\end{proof}

\subsection{Shifted means and random-profile weak comparisons}

Put $T=a^2A_H$, where $A_H$ is the unsigned adjacency matrix, and
$X_\pm=(P^\pm+hI)^{-1}$. We use $h=\kappa_0p^{-4}$ below.
The sharp scalar inequality proved from \eqref{cu:C6}, before choosing
its shift, gives
\[
 t_h^2\le C\{h+p(t_h+b)/(dh)+p^5(t_h+b)/d
                    +p^2/d+p^9/d^2+d^{-1}+e^{-cp}\},
\]
where $t_h=\sup_{i,\pm}(1-\mathbb EX^\pm_{ii}/y_i^\pm)_+$ and
$b=\varepsilon+Cp^4/d+d^{-1}$. For the chosen $p=\lfloor c_0d^{2/17}\rfloor$ and
$h=\kappa_0p^{-4}$, the coefficient of $t_h$ is $o(\sqrt h)$,
and the other constant terms are $o(h)$ as $d\to\infty$ with
$c_0,\kappa_0$ fixed. Thus $t_h\le C\sqrt h$ with an absolute
constant $C$, and
\[
 0\le\mathbb E(G^\pm_{ii}-X^\pm_{ii})/y_i^\pm\le C\sqrt h.
 \tag{S1}\label{eq:mean-shift}
\]
The nonnegative difference is bounded by $G^\pm_{ii}$.
High-moment interpolation therefore preserves this estimate after
multiplication by any fixed-degree product of normalized inverse
diagonals. The interpolation loss is bounded because all parameter
scales are polynomial and $p/\log d\to\infty$.

\begin{lemma}[Random-profile weak loop]\label{lem:random-weak}
Let $X_\epsilon=(P^\epsilon+hI)^{-1}$ and
$X_\nu=(P^\nu+hI)^{-1}$, with $\epsilon,\nu\in\{+1,-1\}$.
Put
\[
 K_\sigma=(X_\epsilon\circ X_\nu)/4,\qquad
 M_\sigma=D_{(X_\epsilon)_{ii}(X_\nu)_{ii}/4},\qquad
 u=1_N/\sqrt d,\quad b_\sigma=4TM_\sigma u,\quad
 w_\sigma=u-\epsilon\nu b_\sigma.
\]
Let $H$ be $1$, $(G^+_{vv}/2)^2$, or $G^+_{vv}G^-_{vv}/4$.
For either $f=u$ or $f=b_\sigma$, set $U=X_\epsilon D_fX_\nu$ and
$F_{ji}=(X_\epsilon)_{ii}U_{ji}-(X_\epsilon)_{ji}U_{ii}$.
Let $D_*$ be one plus the largest physical inverse diagonal of either
branch on $N\cup N_H(N)$ and at $v$. Then
\[
 \left|\mathbb E[H w_\sigma^TK_\sigma f]
                  -\mathbb E[H u^TM_\sigma f]\right|
 \le \mathsf S_f+C\{p/(dh)+1/(dh^{3/2})\},
 \tag{W1}\label{eq:random-weak-error}
\]
where the score retains its physical weight:
\[
 \mathsf S_f=Ca\sum_{i\in N,j\sim i}u_i\mathbb E\bigl[
 |F_{ji}|\{2paH(|G^+_{ij}|+|G^-_{ij}|)+|\partial_{ij}H|\}\bigr].
\]
All its moments below are finite; no a priori smallness of this score
is assumed.
\end{lemma}
\begin{proof}
The following identity holds before taking expectation. For any
possibly random diagonal mask $f$, let $U=X_\epsilon D_fX_\nu$ and
\[
 F_{ji}=(X_\epsilon)_{ii}U_{ji}-(X_\epsilon)_{ji}U_{ii}.
\]
Subtract the scalar root equation multiplied by $U_{ii}$ from the
masked equation multiplied by $(X_\epsilon)_{ii}$ to get
\[
 U_{ii}-f_i(X_\epsilon)_{ii}(X_\nu)_{ii}
       =-\epsilon a\sum_{j\sim i}\sigma_{ij}F_{ji}.       \tag{W2}
\]
All expressions are physical products, with no inverse source or
inverse diagonal. With $f$ temporarily fixed, differentiation in
edge $ij$ gives exactly
\[
 \partial_{ij}F_{ji}=-a\{2\epsilon(X_\epsilon)_{ij}F_{ji}
 +\nu(X_\nu)_{ij}F_{ji}
 +\nu(X_\epsilon)_{ii}(X_\nu)_{ii}U_{jj}
 -\nu(X_\nu)_{ii}(X_\epsilon)_{ij}U_{ij}\}.              \tag{W3}
\]
First-order endpoint integration by parts in (W2), multiplying by
$Hu_i$ and summing $i\in N$, therefore gives
\[
 \mathbb E[H w_\sigma^TK_\sigma f]
       =\mathbb E[H u^TM_\sigma f]+\mathrm{Err}.           \tag{W4}
\]
The leading term in (W3) used here is
$\epsilon\nu a^2(X_\epsilon)_{ii}(X_\nu)_{ii}U_{jj}$.
The determinant derivative is $2pa(G^+_{ij}-G^-_{ij})$.
A derivative of either physical root weight is a bounded-degree
root diagonal product times $aG^\tau_{vi}G^\tau_{vj}$.
Taking absolute values bounds these contributions by $\mathsf S_f$.

Both masks $f=u$ and $f=b_\sigma$ are supported on the deterministic
set $N\cup N_H(N)$, of size at most $d+d^2$.
For the maximum $D_*$ defined above and $k\asymp p$, its
$L^k$ norm is bounded: the factor $(Cd^2)^{1/k}$ is bounded because
$p/\log d\to\infty$. In particular
$\|f\|_\infty\le CD_*^2/\sqrt d$ pointwise.
Put $r_i^\tau=(X_\tau^2)_{ii}$.
Since $X_\tau^2\preceq h^{-1}(G^\tau-X_\tau)$, the mean shift gives
\[
 \mathbb E[D_*^C r_i^\tau]\le Ch^{-1/2}.
 \tag{W5}
\]
For the weighted assertion interpolate $r_i^\tau/y_i^\tau$,
whose pointwise envelope is $h^{-1}h_i^\tau$; the loss
$h^{-O(1/p)}$ is bounded and zero sources follow by continuity.
Write $m=\|f\|_\infty$. Operator norm and Cauchy--Schwarz give
\[
 \|Ue_i\|\le mh^{-1}\sqrt{r_i^\nu},\quad
 |U_{ii}|\le m\sqrt{r_i^\epsilon r_i^\nu},\quad
 \|F_{\cdot i}\|\le2mh^{-1}D_*\sqrt{r_i^\nu}.
\]
Consequently, for $\tau=\epsilon,\nu$,
\[
 \sum_{j\sim i}|(X_\tau)_{ij}F_{ji}|
 \le2mh^{-1}D_*\sqrt{r_i^\tau r_i^\nu},\qquad
 \sum_{j\sim i}|(X_\epsilon)_{ij}U_{ij}|
 \le mh^{-1}r_i^\epsilon.
\]
Using $m\le CD_*^2/\sqrt d$ and (W5), each expectation is at most
$Cd^{-1/2}h^{-3/2}$, including fixed-degree diagonal and root weights.
Multiplication by $a^2u_i$ and summation over $i\in N$ show that the
three secondary terms in (W3) cost $C/(dh^{3/2})$.

For the random mask $f=b_\sigma$, its first derivative is
\[
 \partial_{ij}f_k
 =-2a\sum_{l\in N}T_{kl}u_l
 \{\epsilon(X_\nu)_{ll}(X_\epsilon)_{li}(X_\epsilon)_{lj}
    +\nu(X_\epsilon)_{ll}(X_\nu)_{li}(X_\nu)_{lj}\}.
\]
Using $T_{kl}\le C/d$, $u_l\le d^{-1/2}$ and Cauchy--Schwarz
in $l$ retains the quadratic diagonals:
\[
 \|\partial_{ij}f\|_\infty
 \le\frac{CaD_*}{d\sqrt d}
          \sum_{\tau\in\{\epsilon,\nu\}}\sqrt{r_i^\tau r_j^\tau}.
\]
For any mask $g$, $|F_{ji}(g)|\le Ch^{-1}D_*^2\|g\|_\infty$.
Weighted Cauchy--Schwarz and (W5) therefore imply
\[
 \mathbb E[D_*^C|F_{ji}(\partial_{ij}f)|]
       \le Ca/(d\sqrt d\,h^{3/2}).
\]
There are at most $d^2$ pairs $(i,j)$; the outer factor $a u_i$
thus makes the entire random-mask derivative cost $C/(dh^{3/2})$.

Finally the endpoint remainder is $Cp/(dh)$. Use the uniform incident
row bound with $\rho_{\rm row}=\delta_{\rm row}\le Cp^{-2}$.
For an edge $ij$ let $R_{ij}=|G^+_{ij}|+|G^-_{ij}|$ and
$\ell=\log W$. Exact inverse differentiation gives
\[
 |\ell'|\le CpaR_{ij},\quad
 |\ell''|\le Cpa^2D_*^2,\quad
 |\ell'''|\le Cpa^3R_{ij}D_*^2.
\]
Consequently
\[
 \begin{split}
 |W''|/W&\le Ca^2(p^2R_{ij}^2+pD_*^2),\\
 |W'''|/W&\le Ca^3(p^3R_{ij}^3+p^2R_{ij}D_*^2).
 \end{split}
\]
Differentiating $U$ preserves one $U$ block, or produces
$X_\epsilon D_{\partial f}X_\nu$; further derivatives have the
same form. For $0\le r\le3$,
$\|\partial_{ij}^r f\|_\infty\le C_ra^rD_*^{C_r}/\sqrt d$,
and all extra inverse entries are bounded by physical diagonal products.
Thus for either mask the fixed derivatives of the observable satisfy
\[
 |\partial_{ij}^r(HF_{ji})|
 \le C_r a^r D_*^{C_r}/(\sqrt d\,h),\qquad 0\le r\le3.
\]
The bound includes derivatives of the random mask and physical root
weight, and introduces neither a power of $p$ nor an extra $h^{-1}$. The product rule therefore gives
\[
 W^{-1}|\partial_{ij}^3(WHF_{ji})|
 \le \frac{Ca^3D_*^C}{\sqrt d\,h}
 \{p^3R_{ij}^3+p^2R_{ij}+p^2R_{ij}^2+p+1\}.
\]
Average jointly over $i\in N$, padded to $d$ slots, and a uniform
neighbor $j$, also padded to $d$ slots. Capped moments and
interpolation give $\mathbb E[D_*^C R_{ij}^k]\le C\rho_{\rm row}$
for $k=2,3$, and at most $C\sqrt{\rho_{\rm row}}$ for $k=1$.
Thus the averaged coefficient is
\[
 C\{p^3\rho_{\rm row}+p^2\sqrt{\rho_{\rm row}}+p\}\le Cp.
\]
The first-order endpoint identity, summed with the outer $a u_i$,
therefore has remainder $Cp/(dh)$ rather than $Cp^3/(dh)$.

For completeness this estimate can be transferred along each regular
two-point edge fiber. If the endpoint diagonal sum is at most
$(64pa)^{-1}$, inverse order bounds every diagonal along the segment
by twice its endpoint value. Moreover inverse differentiation gives
\[
 \sup_{s\in[-1,1]}R_{ij}(s)
 \le R_{ij}(\sigma)+CaD_*^2.
\]
Thus one may replace $\rho_{\rm row}$ above by
$\rho_{\rm row}+d^{-1}$, which remains $O(p^{-2})$ in the current
parameter range. Weight ratios on a regular segment are bounded.
On every omitted good endpoint the local diagonal sum exceeds
$c/(pa)$. Using eight extra capped moments bounds its total
contribution, including the first derivative term, by
$Cp^8d^{-7/2}h^{-1}=o(p/(dh))$ for $p\le d^{1/7}$.
All bounds involve physical entries and are uniform at zero sources.

Here $\rho_{\rm row}=\delta_{\rm row}=C\delta/p$ and
$\delta=C\{\varepsilon+p^4/d+(p/d)^{1/3}\}$, so
$\rho_{\rm row}\le Cp^{-2}$ for $p\le d^{1/7}$.
Thus the weak remainder is
\[
 B_0=C\{p/(dh)+1/(dh^{3/2})+e^{-cp}\}.
\]
This proves (W4) with error (W1).
This applies to both stated masks; symmetry of $K_\sigma$ gives
the corresponding first and second weak identities.
\end{proof}

\paragraph{Weighted fresh-star quadratic domination.}
Freeze the graph obtained by deleting a vertex $i$, and let $\xi$ be
the at most $d$ incident signs.  Let $\Phi=\alpha_+^p\alpha_-^p$
be its paired insertion weight, with the conventions and actual-core
normalization of (E4).  Write $G_\tau=(P_\tau)^{-1}$,
$X_\tau=(P_\tau+hI)^{-1}$, and
$C_\tau=(P_{\tau,-i}+hI)^{-1}$.
Here $d^{-C_h}\le h\le1$, $p/\log d\to\infty$, and $p\le d^{1/7}$.

The weights below are fixed-degree products of physical inverse
diagonals, shifted or unshifted. No maximum of inverse diagonals
is included in the Gaussian weight; maxima are used only afterward
in capped-moment interpolation.

For the cavity application set $J=N(i)$, $N=N(v)$,
$u_k=1_N(k)/\sqrt d$, and let $\epsilon\in\{+,-\}$.  Restrict all
matrices below to the deleted core, omitting the zero $i$ coordinate.
The two admissible linear maps from $\mathbb R^J$ to $\mathbb R^J$ are
\[
 \begin{split}
 (T_{\rm dir})_{jl}
   &=1_{\{l\in N\}}(C_\epsilon)_{jl}u_l(C_+)_{ll},\\
 (T_{\rm cav})_{jl}
   &=\sum_{k\in N\setminus\{i,l\}}
               (C_\epsilon)_{jk}u_k(C_+)_{kl}.
 \end{split}                                                    \tag{WT1}
\]
Let $L=T^TT\succeq0$ for either map, or for a sum of a bounded number
of their positive quadratic forms.  Let
\[
 H=H_0(X_\epsilon)_{ii}^2(X_+)_{ii}^2,
\]
where $H_0$ is such a fixed-degree diagonal product whose indices
lie in the radius-two ball under consideration.
For every fixed $M>0$, writing $\vartheta=d^{-M}$, one has
\[
 \boxed{\quad
 \mathbb E[Hq_L(\xi)]\le C\mathbb E[H\operatorname{tr}L]
                                  +C\vartheta+e^{-cp}.
 \quad}                                                         \tag{WT2}
\]
Constants may depend on the fixed degrees, $M$, and $C_h$; they are
uniform in the graph order and in all sources, including limits of
zero sources.  This statement is restricted to the cavity kernels
(WT1), rather than to arbitrary PSD kernels with uncontrolled source
or residual-support dependence.

\paragraph{The weighted Gaussian step.}
For a branch suppress its sign and write $m,n$ for the total powers
of unshifted and shifted denominator determinants in $H$.  The
corresponding part of $\det(P)^pH$ is a product of positive powers of
principal determinants and
\[
 \det(P)^{p-m-n}
            \left(\frac{\det P}{\det(P+hI)}\right)^n.
\]
It is log-concave on $P\succ0$: principal log determinants are
concave, and
\[
 \log\det P-\log\det(P+hI)
       =-\int_0^h\operatorname{tr}(P+tI)^{-1}\,dt
\]
is concave because trace inverse is convex.  The exponent $p-m-n$
is nonnegative.  Extension by zero preserves log-concavity.  The
vertex gauge flip at $i$ takes $\xi$ to $-\xi$ and leaves every
principal determinant unchanged.  Thus the normalized Gaussian law
with density $H\Phi$ is even and has covariance at most the identity
by the Brascamp--Lieb covariance inequality (proved below in the
shifted comparison).  In particular, conditionally on every core,
\[
 \mathsf G[Hq_L\Phi]\le
                   (\operatorname{tr}L)\mathsf G[H\Phi].          \tag{WT3}
\]
This applies to non-core-constant root diagonal weights as well.

\paragraph{Retained comparison terms preserve the unknown mass.}
Put
\[
 x=\mathbb E[Hq_L],\qquad y=\mathbb E[H\operatorname{tr}L],
\]
and let $x_G,y_G$ be their Gaussian numerators divided by the same
full sign insertion normalizer.  Apply (E4) to $Hq_L\Phi$ and
$H(\operatorname{tr}L)\Phi$.  A retained grade $g$ with $l$ active
coordinates has derivative order $n=2(g+l)\le4g$.  At sign endpoints,
derivatives of $H\Phi$, divided by $\Phi$, preserve $H$ up to
fixed-degree diagonal multipliers and coefficient mass $(Cp)^n a^n$.
For inverse diagonals this follows by inverse differentiation and
Cauchy--Schwarz in the associated positive inverse matrix;
there is no inverse-source loss.

At most two derivatives hit $q_L$.  Summing the distinguished
coordinates before estimating them gives the following four bounds,
with the remaining diagonal multipliers suppressed:
\[
 \begin{array}{c|c}
 \text{derivatives on }q_L&\text{sum at grade }g\\ \hline
 0&a^n d^lq_L\le C^g d^{-g}q_L\\
 1&a^{n-1}d^{l-1}\sum_s|(L\xi)_s|
                  \le C^g d^{-g}\sqrt{(\operatorname{tr}L)q_L}\\
 2\text{ on one coordinate}
      &a^{n-2}d^{l-1}\sum_sL_{ss}
                  \le C^g d^{-g}\operatorname{tr}L\\
 2\text{ on distinct coordinates}
      &a^{n-2}d^{l-2}\sum_{s,t}|L_{st}|
                  \le C^g d^{-g}\operatorname{tr}L.
 \end{array}                                                     \tag{WT4}
\]
Here $|(L\xi)_s|^2\le L_{ss}q_L$,
$\sum_s\sqrt{L_{ss}}\le\sqrt{d\operatorname{tr}L}$, and
$\sum_{s,t}|L_{st}|\le d\operatorname{tr}L$.
The last line only occurs when $l\ge2$.

The existing capped-moment interpolation, with the deterministic floor
$\vartheta$, bounds each suppressed fixed-degree multiplier against
the nonnegative masses $Hq_L$ and $H\operatorname{tr}L$.  For example,
the mixed term is bounded by $C\sqrt{(x+\vartheta)(y+\vartheta)}$.
At grade $g$ the extra degree is $O(g)$, and the polynomial-floor
factor is at most $C^g$ since $p/\log d\to\infty$.
Consequently the retained comparison terms sum to
\[
 |x_G-x|+|y_G-y|
 \le C\frac{p^4}{d}(x+y+\vartheta)+e^{-cp}.                       \tag{WT5}
\]

\paragraph{Why the absolute remainder has no hidden source or shift loss.}
For completeness, (WT5) uses the following specific envelope of (WT1).
Let $Q=(P_{+,-i})^{-1}$ be the unshifted plus core.  On residual
support,
\[
 \xi^TQ[J,J]\xi\le Z_i^+/a^2,\qquad
 C_+^2\preceq h^{-1}C_+\preceq h^{-1}Q,
 \qquad
 \xi_l^2\le Z_l^+\,\xi^TQ[J,J]\xi.
\]
The last inequality follows by padding $\xi$ by zero and using
Cauchy--Schwarz in $Q$; $(Q^{-1})_{ll}=Z_l^+$.
These inequalities bound the full and diagonal-subtracted maps in
(WT1) by a fixed polynomial in $d,h^{-1}$, capped core diagonal
factors (including $Z_l^+Q_{ll}$), and at most one factor $Z_i^+$.
The latter cancels against $(X_+)_{ii}^2$ in $H$:
\[
 (X_+)_{ii}\le G^+_{ii}=(Z_i^+\alpha_+)^{-1},
 \qquad (Z_i^+)^{-1}\le C.
\]
Other physical inverse diagonals are bounded on residual support by
their core diagonal times a fixed inverse power of $\alpha_+$ or
$\alpha_-$.  Shifted inverses are bounded by unshifted ones.
Their derivatives cost $a$ per derivative and further residual powers;
they do not cost a new factor $h^{-1}$ per derivative.  Two derivatives
of $q_L$ remove its $\xi$ dependence, and the preceding envelope also
bounds its first derivative by PSD Cauchy--Schwarz.

Thus the grade-$k_*+1$ absolute remainder is at most a fixed polynomial
in $d,h^{-1}$ times
\[
 40^p(Cp)^{4k_*+4}d^{-k_*-1},\qquad
 k_*=\lceil16p/\log d\rceil.
\]
The fixed polynomial is harmless in the exponential reserve of (E6).
This verifies the remainder used in (WT5) without any additional
small-source lower bound.  Finally (WT3), (WT5), and $p^4/d=o(1)$
give $x\le y+C(p^4/d)(x+y+\vartheta)+e^{-cp}$; absorb $x$ to prove
(WT2).

\subsection{Comparison of fresh-star row energies}

\begin{lemma}[Contact row comparison]\label{cu:star-comparison}
At a plus mean contact $\mathbb Eh_v^+=r$, let $\delta$ be the global
concentration envelope in (C2), $N=N(v)$ and
$S_+=\sum_{i\in N}\mathbb E(G^+_{vi})^2$.
Suppose $a^2S_++\mathbb E\operatorname{tr}A^2\le\lambda$,
where $\lambda$ is a polynomial envelope with $\lambda\le C\delta$. Put
\[
 \mathcal R_v=(p-1)\sum_N(G^+_{vi})^2-p\sum_NG^+_{vi}G^-_{vi},
 \qquad \mathcal R=\mathbb E\mathcal R_v.
\]
For the numerator $F$ in \eqref{cu:E5},
\[
 \left|\mathcal R-\frac{\mathbb E_{\nu_K}\mathsf GF}{a^2F_H}\right|
 \le E_{\rm row},
\]
where
\[
 E_{\rm row}
 =C\{p^5\sqrt{\lambda\delta}+p^{13}/d^2+e^{-cp}\}
 \tag{CR1}
\]
\end{lemma}
\begin{proof}
For each free row index $j\in N$ and each incident index $k\in N$, put
\[
 x_j=G^+_{vj},\quad c_{kj}=G^+_{vv}G^+_{kj},\qquad
 z_j=G^-_{vj},\quad d_{kj}=G^-_{vv}G^-_{kj}.
\]
For $D_i=a^{-1}\partial_i$, direct differentiation gives
\[
 \begin{split}
 D_ix_j&=-x_ix_j-c_{ij},\\
 D_ic_{kj}&=-2x_ic_{kj}-x_kc_{ij}-G^+_{vv}G^+_{ki}x_j,
 \end{split}
 \qquad
 \begin{split}
 D_iz_j&=z_iz_j+d_{ij},\\
 D_id_{kj}&=2z_id_{kj}+z_kd_{ij}+G^-_{vv}G^-_{ki}z_j.
 \end{split}
 \tag{CR2}
\]
Thus arbitrary derivative words preserve a plus marked factor in
$\{x_j,c_{kj}\}$ or a minus marked factor in $\{z_j,d_{kj}\}$;
the second index remains the original free index $j$, and every first
index $k$ introduced into a two-index mark is an active derivative
coordinate. Multiplication by the score $2p(x_i-z_i)$ preserves all
existing marks. Starting from $(p-1)x_j^2-px_jz_j$, every monomial
therefore retains two plus marks or one mark of each branch.


Let $I,J$ be independent uniform indices in $N$, padded to $d$ slots
with zeros. The root marked factors satisfy
\[
 \mathbb E|x_J|^2\le C\lambda,\qquad
 \mathbb E|z_J|^2\le C\delta.
\]
For the other marked factors, Schur decomposition gives
\[
 a^4(G^+_{vv})^2\|G^+[N,N]\|_F^2
 \le2(Z_v^+G^+_{vv})^2\operatorname{tr}A^2
       +2\left(a^2\sum_N(G^+_{vi})^2\right)^2.
\]
Both nonnegative random quantities on the right have first moment
$O(\lambda)$ by the hypothesis, before multiplication or squaring, and have
bounded $L^{cp}$ norms by the source cap. High-moment interpolation
therefore yields
\[
 \mathbb E|c_{IJ}|^2\le C\lambda.
\]
The same argument using (C2) gives $\mathbb E|d_{IJ}|^2\le C\delta$.
Here the expectation includes both independent indices $I,J$.
All root source factors remain in the normalized products
$Z_v^\pm G^\pm_{vv}$; no positive lower bound on a source is used.

Only the first two retained grades need the marked bounds.
Their derivative words are $D_i^4$, $D_i^6$ and $D_i^4D_k^4$.
Every active coordinate has a uniform marginal independent of the
free index $J$; hence the preceding moment bounds also apply when
coordinates repeat or $J$ equals an active coordinate.
High-moment interpolation bounds a mixed marked monomial by
$C\sqrt{\lambda\delta}$, and a two-plus monomial by $C\lambda$.
The first two grades consequently cost
$Cp^5\sqrt{\lambda\delta}$ and
$Cp^9d^{-1}\sqrt{\lambda\delta}$.
All grades at least three may use the original unmarked bounds:
they sum geometrically to $Cp^{13}/d^2$.
The final insertion-floor remainder is exponentially small.
This proves (CR1), since $p^4/d=o(1)$.

\end{proof}

\paragraph{Separating incident-row and core marks.}
Suppose the incident-row envelopes are $\lambda_r,\rho_r$, and the
marked-core envelopes are $\lambda_c,\delta_c$, respectively, with
$\lambda_r\le\rho_r$, $\lambda_c\le\delta_c$,
$\lambda_r\le\lambda_c$ and $\rho_r\le\delta_c$.
The row-observable comparison improves to
\[
 E_{\rm row}=C\{p^5\sqrt{\lambda_r\rho_r}
             +p^4\sqrt{\lambda_c\delta_c}+p^{13}/d^2+e^{-cp}\}.
 \tag{CR3}\label{eq:split-row-comparison}
\]
Indeed, in
$\partial_i^4(\Phi\mathcal R_v)=\Phi^{(4)}\mathcal R_v
+\sum_{k=1}^4\binom{4}{k}\Phi^{(4-k)}\partial_i^k\mathcal R_v$,
the first term retains the original free incident factors
$x_j^2$ or $x_jz_j$ and has coefficient $O(p^5a^4)$.
Every other term differentiates the observable at least once,
so its coefficient is at most $Cp^4a^4$ and the core-mark estimate
applies. The marked second grade costs
$Cp^9d^{-1}\sqrt{\lambda_c\delta_c}$, absorbed into the core-mark
term since $p^5/d=o(1)$. The higher-grade tail is unchanged.

\begin{lemma}[Multiplicative shifted trace comparison]\label{cu:shifted-star}
Let $Y_\pm=(P^\pm_{H-v}+hI)^{-1}_{NN}$ and let $E_{\rm row}$
be any available raw comparison error for the row observable $F$.
Put
\[
 \begin{split}
 Q=a^2\mathbb E\{&(p-1)(G^+_{vv})^2\operatorname{tr}Y_+^2
       +pG^+_{vv}G^-_{vv}\operatorname{tr}(Y_+Y_-)\}\ge0,\\
 \eta_{\rm BL}&=C\{p/(dh)+p^4/d\},\qquad
 \epsilon_{\rm BL}=C(p^5+p^2/h)d^{-10}+e^{-cp}.
 \end{split}
\]
Then
\[
 \mathcal R\ge(1-\eta_{\rm BL})Q-E_{\rm row}-\epsilon_{\rm BL}.
 \tag{B1}\label{eq:BLstar}
\]
At the parameters used below, $\eta_{\rm BL}=o(1)$ and
$\epsilon_{\rm BL}=o(p^{-1})$. In particular
$\mathcal R\ge-E_{\rm row}-\epsilon_{\rm BL}$, and
$\mathcal R,E_{\rm row}\le Cp$ imply $Q\le Cp$.
\end{lemma}
\begin{proof}
Fix the core and set
\[
 M=a^2Y_+/Z_v^+,\qquad N=a^2Y_-/Z_v^-,\qquad
 H=2(p-1)M/\alpha+2pN/\beta.
\]
Thus $0\preceq M\preceq A$, $0\preceq N\preceq B$, and
$\|M\|,\|N\|\le Ca^2/h$ because $1/Z_v^\pm\le y_v^\pm\le s$.
Let $\nu$ be the normalized Gaussian law with weight
$\alpha^{p-1}\beta^p$. Its potential, apart from $|x|^2/2$, is
$V=-(p-1)\log\alpha-p\log\beta$, and $\nabla^2V\succeq H$.
The Brascamp--Lieb variance inequality~\cite{BL} for linear functions gives
\[
 \Sigma_\nu\preceq\mathbb E_\nu(I+\nabla^2V)^{-1}
 \preceq\mathbb E_\nu(I+H)^{-1}.
\]
The density is even, so $\Sigma_\nu=\mathbb E_\nu xx^T$.
For completeness, the variance inequality follows by solving
$Lw=f-\mathbb E_\nu f$, where $L=-\Delta+\nabla U\cdot\nabla$ and
$U=|x|^2/2+V$. Integration by parts and weighted Cauchy--Schwarz give
\[
 (\operatorname{Var}_\nu f)^2
 \le \mathbb E_\nu[\nabla f^T(\nabla^2U)^{-1}\nabla f]
       \mathbb E_\nu[\nabla w^T\nabla^2U\nabla w].
\]
The Bochner identity bounds the last factor by
$\mathbb E_\nu(Lw)^2=\operatorname{Var}_\nu f$.
Smooth convex approximation proves the statement for these clipped
powers on their convex support. Since $I-\Sigma_\nu$
is PSD, $A\succeq M$, and $H(I+H)^{-1}\succeq H-H^2$,
\[
 \operatorname{tr}A-\mathbb E_\nu q_A
 \ge\mathbb E_\nu\operatorname{tr}\{M(H-H^2)\}.
\]
Writing $F=(p-1)q_{A^2}\alpha^{p-2}\beta^p
 +pq_{AB}\alpha^{p-1}\beta^{p-1}$, Gaussian integration by parts gives
$2\mathsf GF=\mathsf G[(\operatorname{tr}A-q_A)\alpha^{p-1}\beta^p]$.
Therefore
\begin{align*}
 \mathsf GF\ge\mathsf G\biggl[\Phi\biggl\{
 &(p-1)\frac{\operatorname{tr}M^2}{\alpha^2}
  +p\frac{\operatorname{tr}MN}{\alpha\beta}
  -\frac{\operatorname{tr}(MH^2)}{2\alpha}\biggr\}\biggr].       \tag{B2}
\end{align*}
Write
\[
 X_1=\operatorname{tr}M^2\alpha^{-2},\quad
 X_2=\operatorname{tr}(MN)\alpha^{-1}\beta^{-1},\quad
 x=\mathbb E(X_1+X_2),\quad \vartheta=d^{-10}.
\]
No bound on $x$ is assumed. With $t_+=\alpha^{-1}=D_v^+h_v^+$,
$t_-=\beta^{-1}=D_v^-h_v^-$ and $\mu_A=\operatorname{tr}A$,
$\mu_B=\operatorname{tr}B$, one has
\[
 X_1\le\mu_A^2t_+^2,\qquad X_2\le\mu_A\mu_Bt_+t_-.
\]
These envelopes have bounded $L^{cp}$ moment bases by Schur order,
Minkowski and the capped diagonal moments. Interpolation consequently
gives, for any fixed-degree product $Z$ of normalized inverse diagonals,
\[
 \mathbb E[(X_1+X_2)Z]\le C(x+\vartheta).
 \tag{B3}
\]
Indeed the usual bound is $Cx^{1-C/p}$; split at $x=\vartheta$ and
use $\vartheta^{-C/p}=O(1)$. A product of degree $O(j)$ similarly
costs at most $C^{O(j)}(x+\vartheta)\vartheta^{-Cj/p}$.

Apply \eqref{cu:E4} to the four nonnegative numerators
\[
 \operatorname{tr}M^2\alpha^{p-2}\beta^p,\quad
 \operatorname{tr}(MN)\alpha^{p-1}\beta^{p-1},\quad
 \operatorname{tr}M^2\alpha^{p-3}\beta^p,\quad
 \operatorname{tr}(MN)\alpha^{p-1}\beta^{p-2}.
\]
Every derivative retains its core-constant trace factor. Under the
actual sign law the normalized grade-$j$ correction is bounded by
\[
 (Cp^4/d)^j(x+\vartheta)\vartheta^{-Cj/p}.
\]
The extra factor is absorbed into the geometric base, so the retained
corrections total $C(p^4/d)(x+\vartheta)$.
For the last two numerators the additional $t_+$ or $t_-$ is included
in (B3). The final remainder is exponentially small by \eqref{cu:E6}:
the trace coefficients are bounded by $\mu_A^2$ or $\mu_A\mu_B$,
and require only a fixed number of extra core moments.
All excluded endpoint derivatives still vanish.
Writing $\langle\cdot\rangle_G$ for Gaussian numerator expectation,
averaged over the own core and divided by $F_H$, we obtain
\[
 \begin{aligned}
 \left|\langle\Phi((p-1)X_1+pX_2)\rangle_G
       -\mathbb E((p-1)X_1+pX_2)\right|
 &\le Cp^5d^{-1}(x+\vartheta)+e^{-cp},\\
 \langle\Phi(X_1t_++X_2t_-)\rangle_G
 &\le C(x+\vartheta)+e^{-cp}.
 \end{aligned}
 \tag{B4}
\]
Use $(H_1+H_2)^2\preceq2H_1^2+2H_2^2$ and the operator norm only
once to bound the cubic term in (B2):
\[
 \alpha^{-1}\operatorname{tr}(MH^2)
 \le Cp^2a^2h^{-1}(X_1t_++X_2t_-).
 \tag{B5}
\]
The row-observable comparison and (B4)--(B5), after the raw scale,
give
\[
 \mathcal R\ge Q-E_{\rm row}
                  -C(p^5+p^2/h)(x+\vartheta)-e^{-cp},
\]
where the Schur identities identify
$Q=a^{-2}\mathbb E[(p-1)X_1+pX_2]$ with the stated physical traces.
Since $Q\ge(p-1)a^{-2}x$ and $a^2\asymp d^{-1}$, absorbing the
multiple of $x$ into $Q$ proves \eqref{eq:BLstar}.
All residual factors $t_\pm=D_v^\pm h_v^\pm$ are normalized, so this
argument is uniform at zero sources.
\end{proof}


\subsection{The contact estimate and conclusion}

All expectations in this section use the full Gibbs law, including
expectations of functions of a principal core. Put $q=d-1$ and choose
\[
 p=\lfloor c_0d^{2/17}\rfloor,\quad R^2=4q+4/p,\quad a=R^{-1},
 \quad \frac{q\eta_0^2}{1-\eta_0}=\frac4p,
 \quad s=\frac{2-\eta_0}{1-(1-\eta_0)/q},\quad r=1+\eta_0/2.
 \tag{D1}\label{eq:compressed-parameters}
\]
Thus $r-1\asymp(pd)^{-1/2}$, $s=2+O(d^{-1}+\eta_0)$, and
\[
 \bar L=da^2s^2=\frac{dq(1-\eta_0)}{(q-1+\eta_0)^2},
 \qquad d(r-1)(1-r\bar L)=p^{-1}+O(\eta_0+p^{-1}d^{-1}).
 \tag{D2}\label{eq:compressed-geometric-drift}
\]
We rule out a first contact $\mathbb E h_v^+=r$ in the two source
cubes $0\le y_i^\pm\le s$. Sources equal to zero are interpreted by
continuity; the contacting source is positive. Set $N=N(v)$ and
\[
 \ell_i=a^2y_v^+y_i^+,\quad L=\sum_N\ell_i,\quad
 C_+=\sum_N(c_{vi}^+)^2,\quad D_+=Z_v^+y_v^+=1+L-C_+,
 \quad D_-=Z_v^-y_v^-,
\]
\[
 \begin{gathered}
 \delta_i=r-\mathbb E h_i^+\ge0,\quad
 J=\sum_N\ell_i(\mathcal B(y^+)^{-1})_{vi},\quad
 x_i=G^+_{vi},\quad z_i=G^-_{vi},\\
 S=\sum_N\mathbb E x_i^2,\quad T=\sum_N\mathbb E x_iz_i,
 \quad \mathcal R=(p-1)S-pT.
 \end{gathered}
\]
The deterministic source inequalities give $L\le\bar L$,
$J\ge C_+$ and $C_+/a^2=O(1)$.

\paragraph{Exact loop and bootstrap.}
Define the fourth-order term explicitly by
\[
 P_i=G^+_{vv}G^+_{ii},\quad Q_i=G^-_{vv}G^-_{ii},\quad
 U_i=(p-1)x_i-pz_i,\quad
 V_i=(p-1)(P_i+x_i^2)+p(Q_i+z_i^2),
\]
\[
 \begin{gathered}
 W_i=(p-1)(3P_ix_i+x_i^3)-p(3Q_iz_i+z_i^3),\qquad
 Q_*=(a^4/3)\sum_N\mathbb E K_i,\\
 K_i=x_i(8U_i^3-12U_iV_i+4W_i)
          -3(P_i-x_i^2)(4U_i^2-2V_i).
 \end{gathered}
\]
For a direct derivation let $\alpha=1-\xi^TA\xi$ and
$\beta=1-\xi^TB\xi$ on positive endpoints. The Schur formula gives
$\Phi x_i=-a^{-1}(A\xi)_i\alpha^{p-1}\beta^p$,
$(A\xi)_i/\alpha=-ax_i$ and
$A_{ii}/\alpha=a^2(P_i-x_i^2)$.
Three differentiations therefore give
$\Phi^{-1}\partial_i^3(\Phi x_i)=a^3K_i$.
The endpoint identity supplies the term $Q_*$ directly; no fresh
normalization is required. This polynomial has degree four in
inverse entries and coefficients of degree at most three in $p$.
The corrected loop and the source-maximum covariance inequality are
\[
 D_+r=1+a^2\sum_N\mathbb E
 [G^+_{vv}G^+_{ii}-(2p-1)x_i^2+2px_iz_i]
       +Q_*+O(p^5d^{-2}),
\]
\[
 p\operatorname{Cov}(G^+_{vv},G^+_{ii})
 \le\mathbb E x_i^2-r y_v^+y_i^+(\mathcal B(y^+)^{-1})_{vi}.
\]
Substitute the second formula in the first and use
$\sum c_{vi}^+=L-C_+$. This gives the exact scalar inequality
\[
 \begin{split}
 &(r-1)(1-rL)+r\sum_N\ell_i\delta_i
 +(1-p^{-1})a^2S+(r/p)J+2a^2\mathcal R\\
 &\hspace{30mm}\le Q_*+rC_++O(p^5d^{-2}).
 \end{split}\tag{D3}\label{eq:compressed-drift}
\]
The row-free part of $K_i$ is
$6[(p-1)P_i^2+pP_iQ_i]$. Every other monomial has at least two
root-row factors and coefficient $O(p^3)$. Averaging over a uniform
neighbor index, the precontact row bound gives mean square at most
$C\delta$ for either root row. H\"older with capped high moments
therefore bounds every such monomial by
$C\delta^{1-O(1/p)}\le C\delta$.
Since $a^2d=O(1)$, we obtain
$|Q_*|/a^2\le C(p+p^3\delta)\le Cp$ for the parameters below.
Since
\[
 (r-1)(1-rL)=(r-1)(1-r\bar L)+r(r-1)(\bar L-L),
\]
all terms on the left of \eqref{eq:compressed-drift} other than
$2a^2\mathcal R$ are nonnegative. Consequently
$\mathcal R\le C(p+p^5/d)$. If $\mathcal R\ge-b$, the same
inequality, term by term, gives
\[
 S+a^{-2}\sum_N\ell_i\delta_i\le C(1+p+p^5/d+b),\qquad
 \bar L-L\le\frac{Ca^2(1+p+p^5/d+b)}{r-1}.
 \tag{D4}\label{eq:compressed-bootstrap}
\]
No row saturation is assumed in these estimates.

\paragraph{A single contact bootstrap.}
Fix $0<\kappa_0\le1$ and use
\[
 h=\kappa_0p^{-4},\qquad
 \delta=C\{r-1+p^4/d+(p/d)^{1/3}\}
           \le Cc_0^{17/6}p^{-5/2}.
 \tag{D5}\label{eq:compressed-shift}
\]
All little-$o$ estimates below hold as $d\to\infty$ with
$\kappa_0,c_0$ fixed. Define the unknown deterministic mass
\[
 \lambda=a^2S+\mathbb E\operatorname{tr}A^2+\vartheta,
 \qquad \vartheta=d^{-10}.
\]
By (C2), $\vartheta\le\lambda\le C\delta$, so $\lambda$ is a valid
polynomial envelope in Lemma~\ref{cu:star-comparison}. Its error is
\[
 E_\lambda=C\{p^5\sqrt{\lambda\delta}+p^{13}/d^2+e^{-cp}\}.
\]
The exact Gaussian identity $N_G=2\mathsf GF$ and $N_G\ge0$ give
$\mathcal R\ge-E_\lambda$. Thus \eqref{eq:compressed-bootstrap}
and the already proved $\mathcal R\le Cp$ imply
\[
 a^2S\le C(p+E_\lambda)/d,
 \qquad \mathbb E_{\nu_K}N_G/F_H\le C(p+E_\lambda)/d.
\]
For $w=\operatorname{tr}A^2/(1+\operatorname{tr}A)$, the Gaussian
whitening bound is $N_G\ge f_Gw$.
Apply the unknown-mass interpolation and endpoint-transfer argument
of (B3)--(B4) to $w\Phi$: $w$ has bounded $L^{cp}$ moments and
every derivative retains it. With $x=\mathbb Ew$,
\[
 x\le\mathbb E_{\nu_K}(wf_G)/F_H
                   +C(p^4/d)(x+\vartheta)+e^{-cp}
   \le C(p+E_\lambda)/d+C(p^4/d)(x+\vartheta)+e^{-cp}.
\]
Absorb the multiple of $x$ and use the uniform tail estimate
\eqref{cu:core-tail}. Combining the row and core bounds yields
\[
 \lambda\le Cp/d+Cp^5\sqrt{\lambda\delta}/d
                         +Cp^{13}/d^3+C\vartheta.
\]
Young's inequality absorbs the square-root term, giving
\[
 \lambda\le C\{p/d+p^{10}\delta/d^2+p^{13}/d^3+\vartheta\}
       \le Cp/d=:\lambda_1.
 \tag{D6}\label{eq:compressed-row-bootstrap}
\]
The last inequality follows from $p^7\le d$ and the displayed
formula for $\delta$. In particular $E_\lambda\le Cp$.
Returning to \eqref{eq:compressed-bootstrap} gives
\[
 S\le Cp,\qquad \sum_N\ell_i\delta_i\le Cp a^2,\qquad
 \bar L-L\le\varepsilon_s:=Cp^{3/2}d^{-1/2}.
 \tag{D7}\label{eq:compressed-saturation}
\]
Since $\bar L-L=a^2[ds^2-y_v^+\sum_Ny_i^+]$, with all sources at
most $s$, it follows that
\[
 |y_v^+-2|+\left|\frac{|N|}{d}-1\right|
       +\frac1d\sum_N|y_i^+-2|
       \le C(\varepsilon_s+\eta_0+d^{-1}).
 \tag{D8}\label{eq:compressed-profile-saturation}
\]
No shifted trace budget or source saturation was used to close
$\lambda$.

The global incident-row gain gives the minus envelope
$\delta_{\rm row}=C\delta/p$. Retain the actual plus row mass in
\eqref{eq:split-row-comparison}, taking
$\lambda_r=C(S+1)/d$, $\rho_r=\delta_{\rm row}$,
$\lambda_c=\lambda_1=Cp/d$ and $\delta_c=C\delta$.
The required envelope ordering follows from $S\le Cp$ and
$p^2/d\le C\delta$; the added $1$ supplies a polynomial floor.
With $\Gamma=p^4\sqrt{p\delta/d}$, the comparison error is
\[
 E_{\rm row}=C\{\Gamma\sqrt{S+1}+p^{13}/d^2+e^{-cp}\}.
 \tag{D9a}\label{eq:postcontact-row-error}
\]
In particular $E_{\rm row}\le Cp$, and the positive shifted prefactors
are covered by \eqref{eq:BLstar}.

\paragraph{The contact trace budget.}
Put $Y_\pm=(P^\pm_{H-v}+hI)^{-1}_{NN}$,
$\Omega_+=(G^+_{vv}/2)^2$, and
$\Omega_-=G^+_{vv}G^-_{vv}/4$.
The bounds $\mathcal R,E_{\rm row}\le Cp$ and
\eqref{eq:BLstar} give $Q\le Cp$ directly. Equivalently,
\[
 \mathbb E[\Omega_+d^{-1}\operatorname{tr}Y_+^2]
 +\mathbb E[\Omega_-d^{-1}\operatorname{tr}(Y_+Y_-)]\le C.
 \tag{D9b}\label{eq:contact-trace-budget}
\]
Since $G^+_{vv}\ge1/Z_v^+\ge c$ after saturation, this also bounds
$\mathbb E[d^{-1}\operatorname{tr}Y_+^2]$.
For the full/core correction $\Delta_\pm=X_\pm[N,N]-Y_\pm$,
\eqref{cu:C5} gives
\[
 0\le\operatorname{tr}\Delta_\pm
 \le h^{-1}Z_v^\pm(G^\pm_{vv}-(X_\pm)_{vv}).
\]
The normalized mean-shift estimate and $Z_v^\pm y_v^\pm\le C$
therefore give $\mathbb E[\Omega\operatorname{tr}\Delta_\pm]
\le Ch^{-1/2}$ for either fixed-degree root weight $\Omega$.
High-moment interpolation preserves this bound uniformly at zero
sources. Every shifted compression has norm at most $h^{-1}$, so
expanding either trace product costs $C/(dh^{3/2})$ after normalization.
Thus the same budgets hold for the full compressions $X_\pm[N,N]$.

Since $Q\le Cp$, the multiplicative loss in \eqref{eq:BLstar}
is at most $Cp^5/d+Cp^2/(dh)$.
Consequently
\[
 \begin{split}
 \mathcal R\ge a^2\mathbb E\{&
 (p-1)(G^+_{vv})^2\operatorname{tr}X_+[N,N]^2
 +pG^+_{vv}G^-_{vv}\operatorname{tr}(X_+[N,N]X_-[N,N])\}\\
 &-C\{E_{\rm row}+p^5/d+p^2/(dh)+p/(dh^{3/2})\}.
 \end{split}\tag{D10}\label{eq:compressed-positive-trace}
\]
Both physical root weights have been retained.

\paragraph{Retaining both physical root weights.}
Put
\[
 c=G^+_{vv}/2,\quad d_v=G^-_{vv}/2,\quad
 a_i=(X_+)_{ii}/2,\quad b_i=(X_-)_{ii}/2,
 \quad \Omega_+=c^2,\quad\Omega_-=cd_v.
\]
Use the pointwise PSD matrices
\[
 K_{+,\sigma}=X_+\circ X_+/4,\qquad
 K_{-,\sigma}=X_+\circ X_-/4,
 \quad M_{+,\sigma}=\operatorname{diag}(a_i^2),\quad
 M_{-,\sigma}=\operatorname{diag}(a_ib_i).
\]
Let $b_{\pm,\sigma}=4TM_{\pm,\sigma}u$, and set
\[
 \begin{aligned}
 q_+&=\mathbb E[\Omega_+u^TK_{+,\sigma}u],&
 \zeta&=\mathbb E[\Omega_+u^TK_{+,\sigma}b_{+,\sigma}],\\
 z&=\mathbb E[\Omega_+b_{+,\sigma}^TK_{+,\sigma}b_{+,\sigma}],&
 t_+&=\mathbb E[\Omega_+u^TM_{+,\sigma}u],\\
 \alpha&=\mathbb E[\Omega_+u^TM_{+,\sigma}b_{+,\sigma}],&
 q_-&=\mathbb E[\Omega_-u^TK_{-,\sigma}u],\\
 t_-&=\mathbb E[\Omega_-u^TM_{-,\sigma}u],&
 \beta&=\mathbb E[\Omega_-u^TM_{-,\sigma}b_{-,\sigma}],\\
 z_-&=\mathbb E[\Omega_-b_{-,\sigma}^TK_{-,\sigma}b_{-,\sigma}].
 \end{aligned}
\]
The pointwise matrices and weights are PSD/nonnegative, so
$\zeta\ge0$ and $\zeta^2\le q_+z$ after Cauchy--Schwarz in the
weighted direct-sum space. No random mask is pulled out of an expectation.
For this part of the proof suppress the signing subscript, and write
$q=q_+$ and $t=t_+$. The full-compression trace budget gives $q_+\le C$.
Also $q-t\ge0$ and $\zeta-\alpha\ge0$ exactly, since they are
plus off-diagonal Schur energies.

\paragraph{A nonnegative mask is controlled by its own Schur energy.}
Let $X,Y$ be symmetric positive semidefinite matrices with
$\|X\|,\|Y\|\le h^{-1}$. For a nonnegative mask $b$, put
\[
 B=\operatorname{diag}b,\quad U=XBY,\quad
 K=X\circ Y/4,\quad F_{ji}=X_{ii}U_{ji}-X_{ji}U_{ii}.
\]
Let $I$ be any set of column indices, and suppose $X_{ii}\le D$
for $i\in I$. Then the following deterministic inequality holds:
\[
 \boxed{\quad
 \sum_{i\in I,j}F_{ji}^2
       \le16D^2h^{-2}\,b^TKb.
 \quad}                                                   \tag{BM1}
\]
The sum over $j$ may in particular be restricted to incident edges.
Only the physical diagonal maximum on $I$ is needed.

For the first part of $F$, spectral order and trace positivity give
\[
 \begin{split}
 \sum_{i\in I,j}X_{ii}^2U_{ji}^2
 &\le D^2\operatorname{tr}(U^TU)
   =D^2\operatorname{tr}(BX^2BY^2)\\
 &\le D^2h^{-2}\operatorname{tr}(BXBY).
 \end{split}                                               \tag{BM2}
\]
Each trace comparison is valid by congruence and positivity:
$BX^2B\preceq h^{-1}BXB$, then $Y^2\preceq h^{-1}Y$.

For the second part, Cauchy--Schwarz in the $X$ inner product gives
\[
 U_{ii}^2\le X_{ii}(YBXBY)_{ii}.
\]
Using $(X^2)_{ii}\le h^{-1}X_{ii}$ and the positivity of $YBXBY$,
\[
 \begin{split}
 \sum_{i\in I,j}X_{ji}^2U_{ii}^2
 &\le h^{-1}D^2\operatorname{tr}(YBXBY)\\
 &=h^{-1}D^2\operatorname{tr}(BXBY^2)
 \le D^2h^{-2}\operatorname{tr}(BXBY).
 \end{split}                                               \tag{BM3}
\]
Finally $(r-s)^2\le2r^2+2s^2$ and
$\operatorname{tr}(BXBY)=4b^TKb\ge0$ prove (BM1).
The proof permits a random $b$, since every statement is pointwise.

Apply this with $(X,Y,b)=(X_+,X_+,b_+)$ or
$(X_-,X_+,b_-)$. Capped high-moment interpolation, with
$\theta=d^{-2}$, preserves the respective nonnegative energies $z,z_-$
after fixed-degree diagonal multipliers, so
\[
 \sum_{i\in N,j}\mathbb E[D_*^C\Omega_\pm F_{ji}(b_\pm)^2]
                       \le Ch^{-2}(z_\pm+\theta),\qquad z_+=z.
 \tag{BM4}
\]
Indeed each energy has a polynomial envelope and a bounded high-moment
base. If its first moment exceeds $\theta$, interpolation costs at most
$\theta^{-O(1/p)}h^{-O(1/p)}=O(1)$; otherwise it gives $C\theta$.
The mixed weight need not be bounded below.

\paragraph{The fixed mask is controlled by the off-diagonal propagated energy.}
Set $\beta_h=(dh)^{-1}$ and $\theta=d^{-2}$.
For $\epsilon\in\{+,-\}$, $H\in\{1,\Omega_+,\Omega_-\}$ and
$U=X_\epsilon D_uX_+$, define the same Schur column
$F_{ji}=(X_\epsilon)_{ii}U_{ji}-(X_\epsilon)_{ji}U_{ii}$.
Then
\[
 \sum_{i\in N,j\sim i}\mathbb E[H F_{ji}^2]
 \le C h^{-2}\{\zeta-\alpha+\beta_h(q-t)+\theta\}
                         +C h^{-1/2}.
 \tag{FS1}
\]
The estimate also holds with any fixed-degree physical diagonal
multiplier on the left. Source saturation gives $\Omega_+\ge c>0$,
so $H\le D_*^C\Omega_+$; all such multipliers can be interpolated
against the nonnegative energies on the right, with the floor $\theta$.

To prove this, delete $i$ and write
$C_i^\tau=(P^\tau_{-i}+hI)^{-1}$, embedded with zero $i$ coordinate,
and $\xi=(\sigma_{il})_{l\sim i}$. Exact Schur complementation gives
\[
 F_{ji}=-a(X_\epsilon)_{ii}(X_+)_{ii}
 \sum_{\substack{k\in N\setminus\{i\}\\l\sim i}}
 (C_i^\epsilon)_{jk}u_k(C_i^+)_{kl}\xi_l.
 \tag{FS2}
\]
Split the sum into $k=l$ and $k\ne l$. These are exactly the two
core-fixed linear maps in (WT1). Apply (WT2) with weight
$H(X_\epsilon)_{ii}^2(X_+)_{ii}^2$ to each positive quadratic form.
The total polynomial-floor and exponential errors remain smaller
than $\theta$ after the at most $d$ stars and all displayed polynomial
factors, by choosing the fixed exponent in (WT2) sufficiently large.

For the off-diagonal part, $\|C_i^\epsilon\|\le h^{-1}$ gives
\[
 C h^{-2}\mathbb E\left[D_*^C\Omega_+
 \frac{a^2}{d}\sum_{i\in N}(X_+)_{ii}^2
 \sum_{\substack{l\sim i,\ k\in N\setminus\{i,l\}}}
                        (C_i^+)_{kl}^2\right]+C\theta.
 \tag{FS3}
\]
The exact core/full relation is
$(C_i^+)_{kl}=(X_+)_{kl}-(X_+)_{ki}(X_+)_{il}/(X_+)_{ii}$.
Multiply by $(X_+)_{ii}^2$ before estimating its square; this
introduces no inverse source. The first term is bounded by
\[
 \frac{a^2}{d}\sum_{i\in N}(X_+)_{ii}^2
       \sum_{l\sim i,\ k\in N\setminus\{l\}}(X_+)_{kl}^2
       =4u^T(K_+-M_+)b_+.
\]
For the rank-one term, $\sum_l(X_+)_{il}^2\le h^{-1}(X_+)_{ii}$ gives
\[
 \frac{a^2}{d}\sum_{i,k\in N,\ k\ne i}(X_+)_{ki}^2
                    \sum_{l\sim i}(X_+)_{il}^2
 \le\frac{CD_*}{dh}\frac1d
           \sum_{i,k\in N,\ k\ne i}(X_+)_{ki}^2.
\]
The two weighted first moments are respectively
$4(\zeta-\alpha)$ and $4(q-t)$.
Their polynomial envelopes and the capped high moments give
$\mathbb E[D_*^C Y]\le C(\mathbb EY+\theta)$ for either weighted
nonnegative energy $Y$. This proves the first term in (FS1).

For the direct part, (WT2) instead gives
\[
 Ca^2\sum_{i\in N,\ k\in N\cap N(i)}u_k^2
 \mathbb E\left[H(X_\epsilon)_{ii}^2(X_+)_{ii}^2
                (C_i^+)_{kk}^2((C_i^\epsilon)^2)_{kk}\right]
 +C\theta.
\]
Here $(C_i^+)_{kk}\le(X_+)_{kk}$. Do not square this matrix order.
Use the exact Schur subtraction to obtain
\[
 (X_\epsilon)_{ii}^2((C_i^\epsilon)^2)_{kk}
 \le2(X_\epsilon)_{ii}^2(X_\epsilon^2)_{kk}
              +2(X_\epsilon)_{ik}^2(X_\epsilon^2)_{ii}.
\]
The mean Ward bound (W5), with physical diagonal multipliers, bounds
both expectations by $Ch^{-1/2}$. Since
$a^2\sum_{i,k\in N}u_k^2 1_{\{i\sim k\}}\le C$, the entire direct
part costs $Ch^{-1/2}$. This proves (FS1).

\paragraph{Closing the four weak errors.}
The generic weak score is estimated with its physical weight retained:
the determinant contribution is $2paH|F|(|G^+_{ij}|+|G^-_{ij}|)$,
and the root derivative contribution is $|F\partial_{ij}H|$.
For either nonconstant root weight, positivity gives
\[
 \frac{|\partial_{ij}H|^2}{H}
 \le Ca^2D_*^C\{(G^+_{vi})^2+(G^-_{vi})^2\},
 \tag{FS4}
\]
with the quotient interpreted by continuity at $H=0$.
Indeed $(G^\tau_{vi}G^\tau_{vj})^2/G^\tau_{vv}
\le G^\tau_{jj}(G^\tau_{vi})^2$; no inverse source is retained.
Joint-index Cauchy--Schwarz, (GR1), and (FS1) therefore give
\[
 \mathsf S_u\le
 Ce\sqrt{\zeta-\alpha+\beta_h(q-t)+\theta}+Ce h^{3/4},
 \qquad e=\frac{p\sqrt{\delta_{\rm row}}}{\sqrt d\,h}.
 \tag{FS5}
\]
The same calculation with (BM1) gives
$\mathsf S_{b_+}\le Ce\sqrt{z+\theta}$ and
$\mathsf S_{b_-}\le Ce\sqrt{z_-+\theta}$, where
$z=\mathbb E[\Omega_+b_+^TK_+b_+]$ and
$z_-=\mathbb E[\Omega_-b_-^TK_-b_-]$.
No comparison between plus and minus masks is needed.

Let $r_0=q-\zeta-t$ and $r_1=\zeta-z-\alpha$ be the two actual
plus residuals. The weak loop, applied directly to $f=u$ and $f=b_+$,
gives
\[
 \begin{split}
 |r_1|&\le Ce\sqrt{z+\theta}+CB_0,\\
 |r_0|&\le Ce\sqrt{z+|r_1|+\beta_h(q-t)+\theta}
                                      +Ce h^{3/4}+CB_0,\\
 B_0&=p/(dh)+1/(dh^{3/2})+e^{-cp}.
 \end{split}\tag{FS6}
\]
The contact trace budget gives $q\le C$, hence $0\le q-t\le C$.
Moreover $e\sqrt{|r_1|}\le C(e\sqrt z+e^2+B_0+e\sqrt\theta)$.
Thus every plus or mixed weak residual, with either $f=u$ or its
nonnegative mask $f=b_\pm$, is at most
\[
 E_c=C\{e\sqrt z+e\sqrt{z_-}+e^2+B_0
                  +e[h^{3/4}+(dh)^{-1/2}+\sqrt\theta]\}.
 \tag{FS7}
\]
The mixed identities use linearity of $F$ and (BM1); their fixed-mask
part is covered by the same plus energies in (FS5). All comparison
normalizers are the actual full sign normalizer as in (WT2).

Consequently the four identities, each with absolute error $O(E_c)$, are
\[
 q_+-\zeta=t_++O(E_c),\qquad \zeta-z=\alpha+O(E_c),
\]
\[
 \mathbb E[\Omega_-u^TK_-(u+b_-)]=t_-+O(E_c),\qquad
 \mathbb E[\Omega_-b_-^TK_-(u+b_-)]=\beta+O(E_c).
\]
Every random mask remains inside its expectation.

\paragraph{Averaged profile inequalities.}
Let
\[
 \chi_+=\sqrt{\varepsilon/p+\varepsilon^2+p/d}
             +\varepsilon_s+\eta_0+d^{-1},\qquad
 \rho_v=\varepsilon/p+\varepsilon^2,
 \quad \xi=\sqrt h+\chi_++\sqrt{\rho_v}+\varepsilon+d^{-1}.
\]
The $k=2$ case of (F2) gives
$\mathbb E(h_i^+-1)^2\le2\delta_i+C(\varepsilon/p+\varepsilon^2)$,
while source saturation and (D7) give
$\sum_{i\in N}y_i^+\delta_i\le Cp$.
Weighted Cauchy--Schwarz and the source bounds yield the stated
$\chi_+$ estimate. The contact neighbor estimates and the mean shift bound give
$d^{-1}\sum_N\mathbb E|a_i-1|\le C(\sqrt h+\chi_+)$,
also with any fixed-degree diagonal multiplier. At the contacting
root, the variance bound and source saturation give
$\mathbb E|c-1|\le C(\sqrt{\rho_v}+\varepsilon_s+\eta_0+d^{-1})$.
All assertions are uniform at zero sources.
By the fourth moment cap and Holder,
\[
 0\le t_+,t_-\le1+C\varepsilon.
 \tag{RC1}
\]
For example, $\mathbb E[cd_va_ib_i]$ is bounded above by the product
of the four normalized fourth-moment norms times $(s/2)^4$.

Write $P=A_H/d$ and $\alpha_0=u^TPu$. Replacing $4T$ by $P$ costs
$O(d^{-1})$ in these scalar profiles. Replacing all plus factors in
$\alpha$ by one gives $|\alpha-\alpha_0|\le C\xi$.
In $\beta$ replace only $c,a_i,a_j$ by one. The remaining product
satisfies
\[
 \mathbb E[d_v b_i b_j]
 \le(s/2)^3\mathbb E[h_v^-h_i^-h_j^-]\le1+C\varepsilon.
\]
The error is $C\xi$ by the preceding weighted mean estimates.
Consequently
\[
 \beta\le\alpha+C\xi.                                  \tag{RC2}
\]
No lower bound on a minus source, and no pointwise cap on a random
profile, is used.

\paragraph{A reserve in both quadratic mask energies.}
Put $A_p=p-1$. The preceding four weak identities, the PSD kernels,
and (RC1)--(RC2) imply
\[
 \boxed{\quad
 A_p(q_+-t_+)+p(q_--t_-)
 \ge\frac{A_p}2 z+\frac p2 z_-
          -\frac1{4A_p}-Cp(E_c+\xi).
 \quad}                                                    \tag{JR2}
\]
No upper bound on either quadratic energy is required.

Here is a source-uniform perturbation proof. Plus-kernel PSD gives
$\zeta^2\le q_+z$; combining this with the four weak identities and the upper bound
on $t_+$ yields
\[
 \alpha\le\frac\zeta{1+\zeta}+C(E_c+\varepsilon).
\]
This follows by substituting
$q_+\le\zeta+1+C(E_c+\varepsilon)$ and
$z=\zeta-\alpha+O(E_c)$; the resulting rational expression is
uniformly Lipschitz in its additive error. When $q_+=0$, PSD
forces $\zeta=0$ and the same conclusion follows from the four weak identities.
Choose
\[
 \widehat\alpha=(\alpha-C_0(E_c+\varepsilon))_+
\]
with an absolute sufficiently large $C_0$. Then
$0\le\widehat\alpha\le\zeta/(1+\zeta)<1$ and
$|\alpha-\widehat\alpha|\le C_0(E_c+\varepsilon)$.

For the mixed kernel, put $w=u+b_-$ and
$\lambda_m=\beta/(t_-+\beta)\in[0,1]$, with $\lambda_m=0$ when the denominator
vanishes. The PSD test $b_--\lambda_m w$ and the four weak identities give
\[
 z_-\ge2\lambda_m\beta-\lambda_m^2(t_-+\beta)-CE_c
      =\frac{\beta^2}{t_-+\beta}-CE_c
      \ge\frac{\beta^2}{1+\beta}-C(E_c+\varepsilon).
 \tag{JR3}
\]
The error is uniform because $\lambda_m\le1$; no division by a small
$t_-$ is used. Expanding the four weak identities also gives
$q_--t_-=z_--\beta+O(E_c)$.

Write $\mathscr P=A_p(q_+-t_+)+p(q_--t_-)$. The last identities imply
\[
 \mathscr P-\frac{A_p}2z-\frac p2z_-
 \ge\frac{A_p}2\zeta+\frac{A_p}2\alpha+\frac p2z_--p\beta-CpE_c.
\]
Apply (JR3). The function
$f(x)=x^2/[2(1+x)]-x$ is decreasing and $1$-Lipschitz on
$[0,\infty)$. Thus $\beta\le\alpha+C\xi$ allows replacement of
$f(\beta)$ by $f(\alpha)-C\xi$. Replacing $\alpha$ by $\widehat\alpha$
costs at most $Cp(E_c+\varepsilon)$. Finally
$\zeta\ge\widehat\alpha/(1-\widehat\alpha)$ gives
\[
 \begin{split}
 \frac{A_p}2\zeta+\frac{A_p}2\widehat\alpha+p f(\widehat\alpha)
 &\ge-\widehat\alpha+\frac{\widehat\alpha^2}{2}
       \left\{\frac{A_p}{1-\widehat\alpha}+\frac p{1+\widehat\alpha}\right\}\\
 &\ge A_p\widehat\alpha^2-\widehat\alpha\ge-\frac1{4A_p}.
 \end{split}
\]
This proves (JR2), with the unchanged scalar loss
$c_p=1/[4(p-1)]$.

Substitute (FS7) and absorb $Cpe\sqrt z$ and $Cpe\sqrt{z_-}$
into the two positive reserves by Young's inequality. The result is
\[
 \begin{split}
 &(p-1)(q_+-t_+)+p(q_--t_-)\\
 &\quad\ge-\frac1{4(p-1)}
 -Cp\{e^2+B_0+e[h^{3/4}+(dh)^{-1/2}+\sqrt\theta]+\xi\}.
 \end{split}\tag{RC3}
\]

\paragraph{Exact retained fourth-order center.}
Define the unshifted scalar centers
\[
 t_{+,0}=\frac1{16d}\mathbb E[(G^+_{vv})^2\sum_N(G^+_{ii})^2],\qquad
 t_{-,0}=\frac1{16d}\mathbb E[G^+_{vv}G^-_{vv}
                                      \sum_NG^+_{ii}G^-_{ii}],
\]
and $L_d=16da^2=4+O(d^{-1})$.
The mean shift estimate, with fixed-degree physical diagonal weights,
gives
\[
 0\le t_{+,0}-t_+\le C\sqrt h,\qquad
 0\le t_{-,0}-t_-\le C\sqrt h.                         \tag{RC4}
\]
The positive full shifted traces in the BL comparison equal
$L_d[(p-1)q_++pq_-]$. The row-free part of $Q_*/a^2$ equals
\[
 2L_d[(p-1)t_{+,0}+pt_{-,0}]                            \tag{RC5}
\]
exactly. Thus one may center $\mathcal R$ at
$L_d[(p-1)t_{+,0}+pt_{-,0}]$, without replacing either root diagonal.
Only the $O(p\sqrt h)$ shift error remains.

For the row-containing remainder, use the deterministic envelopes
$d^{-1}\sum_N\mathbb Ex_i^2\le\lambda_1=Cp/d$ and
$d^{-1}\sum_N\mathbb Ez_i^2\le\delta_{\rm row}$.
At $x_i=0$, the row-containing remainder is exactly
$-6p(2p-1)P_i z_i^2\le0$. Subtracting this term leaves only monomials
with a plus row factor and at least one further row factor. Hence
\[
 Q_*/a^2\le2L_d[(p-1)t_{+,0}+pt_{-,0}]
                    +Cp^3\sqrt{\lambda_1\delta_{\rm row}}.          \tag{D13}\label{eq:compressed-centering}
\]



Combining the rational inequality, the exact centers and
\eqref{eq:compressed-positive-trace} gives
\[
 \mathcal D:=\mathcal R-L_d[(p-1)t_{+,0}+pt_{-,0}]
      \ge-L_dc_p-\mathcal E,
 \qquad c_p=\frac1{4(p-1)},
 \tag{D12}\label{eq:compressed-residual}
\]
where the error is accounted for below.

\paragraph{The final drift.}
Moreover $C_+/a^2=4+O(\varepsilon_s+\eta_0+d^{-1})$.
For
\[
 \mathfrak D=(2p-1-p^{-1})a^2S-2pa^2T+(r/p)J-Q_*-rC_+,
\]
use $J\ge C_+$ in \eqref{eq:compressed-centering} to obtain
\[
 \mathfrak D/a^2\ge2\mathcal D+(1-p^{-1})(S-4)-\mathcal E.
 \tag{D14}\label{eq:compressed-centered-drift}
\]
The exact row equation
$a\sum_N\sigma_{vi}x_i=1-Z_v^+G^+_{vv}$ gives
\[
 S\ge\frac{(rD_+-1)^2}{a^2|N|}
       \ge4-C(\varepsilon_s+\eta_0+d^{-1}).
 \tag{D15}\label{eq:compressed-row-lower}
\]

Retain the positive row surplus in \eqref{eq:compressed-centered-drift}.
By \eqref{eq:compressed-row-lower} and Young's inequality,
\[
 C\Gamma\sqrt{S+1}\le\tfrac14(S-4)
 +C\{\Gamma+\Gamma^2+\varepsilon_s+\eta_0+d^{-1}\}.
\]
After absorbing the first term, the remaining error is bounded by
an absolute constant times
\[
 \begin{split}
 &\Gamma+\Gamma^2+p^{13}/d^2+p^5/d+p^2/(dh)+p/(dh^{3/2})
                  +p^3\sqrt{\lambda_1\delta_{\rm row}}\\
 &\qquad+p\{B_0+e^2+e[h^{3/4}+(dh)^{-1/2}+\sqrt\theta]+\xi\},
 \end{split}
\]
plus $\epsilon_{\rm BL}=o(p^{-1})$ and exponentially small terms.
At $p=\lfloor c_0d^{2/17}\rfloor$, the critical terms satisfy
\[
 \begin{aligned}
 p\sqrt h&=\sqrt{\kappa_0}\,p^{-1},\\
 \Gamma&\le Cc_0^{17/3}p^{-1},\\
 pe^2&\le C\kappa_0^{-2}c_0^{34/3}p^{-1}.
 \end{aligned}\tag{D16}\label{eq:compressed-ledger}
\]
For fixed $c_0,\kappa_0$, every other term is $o(p^{-1})$. Indeed
\[
 \begin{gathered}
 d\asymp p^{17/2},\quad\delta=O(p^{-5/2}),\quad
 \delta_{\rm row}=O(p^{-7/2}),\quad\lambda_1=O(p^{-15/2}),\\
 \varepsilon=O(p^{-19/4}),\quad\rho_v=O(p^{-23/4}),\quad
 \varepsilon_s=O(p^{-11/4}),\quad \chi_+=O(p^{-11/4}),\quad e=O(p^{-1}),
 \end{gathered}
\]
which gives
\[
 \begin{gathered}
 \Gamma^2=O(p^{-2}),\quad p^{13}/d^2=O(p^{-4}),\quad p^5/d=O(p^{-7/2}),\\
 p^2/(dh)=O(p^{-5/2}),\quad p/(dh^{3/2})=O(p^{-3/2}),\quad
 pB_0=O(p^{-3/2}),\\
 peh^{3/4}=O(p^{-3}),\quad pe(dh)^{-1/2}=O(p^{-9/4}),\quad
 pe\sqrt\theta=O(p^{-17/2}),\\
 p^3\sqrt{\lambda_1\delta_{\rm row}}=O(p^{-5/2}),\quad
 p(\chi_++\sqrt{\rho_v}+\varepsilon)=O(p^{-7/4}).
 \end{gathered}
\]
Thus the total error, including the final drift errors, is at most
\[
 C\{\sqrt{\kappa_0}+c_0^{17/3}
                 +\kappa_0^{-2}c_0^{34/3}\}p^{-1}+o(p^{-1}).
\]
First choose $\kappa_0$ sufficiently small and then $c_0>0$
sufficiently small so that this coefficient is less than $1/4$.
All choices are absolute.
From \eqref{eq:compressed-residual}--\eqref{eq:compressed-row-lower}
and $L_d=16da^2$, $da^2=1/4+O(d^{-1})$ we get
$d\mathfrak D\ge-2c_p-\mathcal E_0$, where the total error
$\mathcal E_0$ satisfies $\mathcal E_0\le(1/4+o(1))p^{-1}$
by the preceding choice of constants. But
\eqref{eq:compressed-drift} says
\[
 (r-1)(1-r\bar L)+r(r-1)(\bar L-L)
       +r\sum_N\ell_i\delta_i
       \le-\mathfrak D+O(p^5d^{-2}).
\]
The last two terms on its left are nonnegative. Multiplying by $d$
and using \eqref{eq:compressed-geometric-drift} gives the contradiction
\[
 p^{-1}\le 2c_p+(1/4+o(1))p^{-1}
                   =(3/4+o(1))p^{-1}.
\]

\paragraph{Closing the induction.}
Induct on graph order at this fixed $d$. The deletion
reparametrization places every proper core in a smaller source
cube; its induction hypothesis and the insertion estimate make
the full normalizer positive before a possible first mean contact.
The exact moment maximum then supplies every moment used above.
Starting from zero sources, continuity and the contradiction rule
out such a contact throughout both cubes. All degree thresholds
are independent of graph order. At $y^+=y^-=s\mathbf1$, the
Bethe diagonal satisfies $Z_i\le1$. Any signing in the nonempty
support has $P^\pm\succ0$, hence
$I\pm aA_\sigma=P^\pm+\operatorname{diag}(1-Z)\succ0$.
Therefore
\[
 \gamma_d^2\le4(d-1)+4/\lfloor c_0d^{2/17}\rfloor
                \le4(d-1)+Cd^{-2/17}.
\]

\section{Lower bound}

\begin{proof}
Put $e=q^{-1/2}$, $t=e^2$, $k=q-1$, $p=k/q=1-t$, and
$L=\lceil q^6\rceil$. We first choose a simple $k$-regular graph $F$
on $3N$ vertices, partitioned into triples, such that
\begin{equation}\label{eq:lower-host}
 \left\|q^{-1}A_F\bigm|_{\mathbf1^\perp}\right\|
 \le\frac{2\sqrt{k-1}+1}{q},\qquad
 \operatorname{girth}(J)>2L+6,\qquad
 \operatorname{MaxCut}(J)\le\frac{DN}{4}
       +\frac{(P_*+\epsilon_D)N\sqrt D}{2},\quad D=3k,
\end{equation}
where $J$ is obtained by contracting the triples and
$\epsilon_D>0$ is a deterministic sequence tending to zero.
To obtain $F$, pair uniformly the $k$ half-edges at each of its
$3N$ vertices. Grouping triples gives the $D$-regular configuration
model on $N$ vertices. For fixed $q,L$, as $N\to\infty$ through
admissible parities, the contracted girth condition has positive
limiting probability by the short-cycle law~\cite{MWW}; it implies
that both $J$ and $F$ are simple. Conditional on simplicity, $F$
is uniform simple regular, so its spectral condition fails with
probability tending to zero by Friedman's theorem~\cite{Bordenave}.
Likewise, conditional on simplicity, $J$ is uniform simple
$D$-regular. The random-regular MaxCut theorem of
Dembo--Montanari--Sen~\cite[Theorem 1.6(b)]{DMS} gives the displayed
cut bound with probability tending to one, for a suitable
$\epsilon_D\to0$. Intersecting these events with the positive-probability
girth event proves that all conditions in \eqref{eq:lower-host}
hold for some finite $N$.
Add a triangle on each triple to obtain $G$. The contracted girth
makes $G$ simple, and the spectral condition makes it connected;
its degree is $q+1$.

Fix $c\in[1,3]$, set $z=2+ce^6$, $\delta=ce^6$, and suppose that
$\|A_\sigma(G)\|<z\sqrt q$. All $O(\cdot)$ constants below are absolute,
uniform in $c$, $N$, and $\sigma$. Write $\tau_v\in\{\pm1\}$ for the
flux of the triangle containing $v$, and $\varepsilon\in\{\pm1\}$.
For each directed external edge $v\to w$, let
$g^{(j)}_{v\to w,\varepsilon}$ be the root diagonal of
$(zI-\varepsilon eA_\sigma)^{-1}$ on the depth-$j$ tree of triangles
on the $v$ side with $vw$ omitted; depth zero is the root triangle.
The girth condition makes these trees and their one-block extensions
through depth $L+1$ induced subgraphs of $G$. Their matrices are
positive definite, so principal-submatrix inversion and insertion
of a message at the neighboring triangle give
\[
 1/z\le g^{(j)}\le g^{(j+1)},\qquad g^{(j)}<z/t\le3q
 \qquad(0\le j\le L).
\]
For the upper bound, positivity of the neighboring effective diagonal
$z-t\sum g^{(j)}$ bounds each summand by $z/t$.
Let $\mathbb E$ average over directed edges of $F$ and both signs
$\varepsilon$; for vertex fields it is also the uniform vertex/sign
average. Some $1\le j\le L$ has
$\mathbb E(g^{(j)}-g^{(j-1)})\le3q/L$. Set
\begin{equation}\label{eq:lower-defect}
 g=g^{(j-1)},\quad \widehat g=g^{(j)},\quad x=g-1,\quad
 h=g^{-1}-\widehat g^{-1};
 \qquad 0\le h\le3,\quad\mathbb Eh\le27q/L=O(e^{10}).
\end{equation}

Define
\[
 (Pf)_{v,\varepsilon}=t\sum_{u\sim_Fv}f_{u\to v,\varepsilon},
 \qquad (Qf)_{v\to w,\varepsilon}=(Pf)_{v,\varepsilon}
                                  -tf_{w\to v,\varepsilon}.
\]
Their row sums are $p$ and $1-2t$, and
$\mathbb EQf=(1-2t)\mathbb Ef$. For a triangle $v,a,b$, let
$a_u=1+t+\delta-(Px)_{u,\varepsilon}$. Switching its internal signs
shows that its self-energy at $v$ is
\begin{equation}\label{eq:lower-recursion}
 S_{v,\varepsilon}
 =t(1,1)\begin{pmatrix}a_a&-\varepsilon\tau_ve\\
                  -\varepsilon\tau_ve&a_b\end{pmatrix}^{-1}\binom11,
 \qquad
 g^{-1}=1+2t+\delta-Qx-S+h.
\end{equation}
The displayed matrix is positive definite. Its inverse quadratic
form is convex in its diagonal entries. Its tangent at
$a_a=a_b=1+t$ therefore gives
\begin{equation}\label{eq:lower-tangent}
 S=s_{\varepsilon\tau}
   -k_{\varepsilon\tau}\{2\delta-(Px)_a-(Px)_b\}+T,\qquad T\ge0,
 \quad
 s_\eta=\frac{2t}{1+t-\eta e},\quad
 k_\eta=\frac{t}{(1+t-\eta e)^2}.
\end{equation}
Write $s_\eta=s_0+\eta s_1$ and $k_\eta=k_0+\eta k_1$. Directly,
\begin{equation}\label{eq:lower-coefficients}
 s_0=2t-2e^6+O(e^8),\quad
 s_1=2e^3+O(e^5),\quad k_0=t+t^2+O(e^6),\quad
 k_1=2e^3+O(e^5),\quad k_0p-t=O(e^6).
\end{equation}
Put
\[
 V=\mathbb E(x^2/g),\qquad M_1=\mathbb E|x|,\qquad
 C_1=\mathbb E_{v,\varepsilon}[\varepsilon\tau_v(Px)_{v,\varepsilon}].
\]
Weighted Cauchy--Schwarz gives
$M_1^2\le V\mathbb Eg\le V(1+M_1)$, hence
$M_1\le V+\sqrt V$, and $|C_1|\le pM_1$.
Since $\tau$ is constant on each triangle, averaging
\eqref{eq:lower-tangent} and \eqref{eq:lower-recursion} gives the exact identity
\begin{equation}\label{eq:lower-energy}
 (1+2k_0)\delta
 =V+(s_0-2t)+2(k_0p-t)\mathbb Ex+2k_1C_1+\mathbb ET-\mathbb Eh.
\end{equation}
Here one uses $\mathbb E(g^{-1})=1-\mathbb Ex+V$.
Discarding $T\ge0$ first yields
$V\le O(e^6)+O(e^3)(V+\sqrt V)$; absorption, and then
\eqref{eq:lower-energy} again, imply
\begin{equation}\label{eq:lower-small}
 V=O(e^6),\qquad M_1=O(e^3),\qquad\mathbb ET=O(e^6),\qquad
 \mathbb E|S-2t-2\varepsilon\tau e^3|=O(e^5).
\end{equation}
The last assertion follows from \eqref{eq:lower-tangent} and
\eqref{eq:lower-coefficients}.

To retain the first correction, set $u=g^{-1}-1$, $v=x^2/g$, so
$x=-u+v$, and let $U_{v,\varepsilon}$ average $u_{v\to w,\varepsilon}$
over its $k$ outgoing edges. Since $u^2\le3x^2/g$ and
\eqref{eq:lower-recursion} has the form
$u_{v\to w}=F_{v,\varepsilon}+tg_{w\to v}+h_{v\to w}$,
\begin{equation}\label{eq:lower-projection}
 \|U\|_2=O(e^3),\qquad
 \mathbb E|u-U|\le2tM_1+2\mathbb Eh=O(e^5).
\end{equation}
The same recursion, using $x+u=v$, gives
\begin{equation}\label{eq:lower-forcing}
 x-Qx=2\varepsilon\tau e^3+r,\qquad\mathbb E|r|=O(e^5).
\end{equation}
Multiplying by $\varepsilon$ and averaging shows that the vertex mean
$\mu$ of $\tau$ satisfies
$2t\mathbb E(\varepsilon x)=2\mu e^3+O(e^5)$, whence $\mu=O(e^2)$.

Let $\mathsf P=q^{-1}A_F$ act on vertex fields, and use uniform vertex
inner products, also averaging $\varepsilon$ when present. Replacing
$x=-u+v$ and using \eqref{eq:lower-projection}, with
$\mathbb Ev=O(e^6)$, gives
\[
 C_1=-\langle\varepsilon\mathsf P\tau,U\rangle+O(e^5),\qquad
 U-\alpha\mathsf P U=-2\varepsilon\tau e^3+R_0,
 \quad \mathbb E|R_0|=O(e^5),\quad \alpha=(k-1)/k.
\]
The second equality is the outgoing-edge average of
\eqref{eq:lower-forcing}; the factor $\alpha$ excludes the reversed edge.
By \eqref{eq:lower-host} and $\mu=O(e^2)$,
\[
 \|\mathsf P^2\tau\|_2
 \le p^2|\mu|+\|\mathsf P|_{\mathbf1^\perp}\|^2\|\tau-\mu\|_2
 =O(e^2),\qquad \kappa:=\langle\tau,\mathsf P\tau\rangle=O(e).
\]
Pair the preceding equation for $U$ with $\varepsilon\mathsf P\tau$.
Its error is $O(e^5)$ since $\|\mathsf P\tau\|_\infty\le p$, and
$|\langle\varepsilon\mathsf P^2\tau,U\rangle|=O(e^5)$. Thus
\begin{equation}\label{eq:lower-correlation}
                         C_1=2e^3\kappa+O(e^5).
\end{equation}

Finally put $B_1=\mathbb E[\varepsilon\tau_vx_{v\to w,\varepsilon}]$.
External regularity gives
$\mathbb E[\varepsilon\tau_v(Qx)_{v\to w,\varepsilon}]=\alpha C_1$,
so \eqref{eq:lower-forcing} implies
$B_1=2e^3+\alpha C_1+O(e^5)$.
Since $B_1^2\le V\mathbb Eg=V(1+O(e^3))$, it follows that
\[
 V\ge4e^6+4e^3C_1-O(e^8).
\]
Insert this into \eqref{eq:lower-energy}, discard $\mathbb ET\ge0$,
and use \eqref{eq:lower-correlation}:
\[
 \delta\ge V-2e^6+4e^3C_1-O(e^8)
          \ge2e^6+16e^6\kappa-O(e^8).
\]
Since $\tau$ is constant on triples, the cut bound gives
\[
 \kappa=\frac{\tau^TA_J\tau}{3Nq}
 =\frac{DN-4\operatorname{Cut}_J(\tau)}{3Nq}
 \ge-\frac{2(P_*+\epsilon_D)\sqrt D}{3q}
 \ge-\frac{2(P_*+\epsilon_D)}{\sqrt3}\,e.
\]
Consequently the supposed signing requires
$c\ge2-32(P_*+\epsilon_D)e/\sqrt3-C_0e^2$, for an absolute
$C_0\ge1$. Choose
$c=2-32(P_*+\epsilon_D)e/\sqrt3-2C_0e^2\in[1,3]$ for sufficiently
large $q$. This contradicts the requirement. The finite graph was
chosen independently of $c$ and $\sigma$. Since $\epsilon_D\to0$,
multiplying $z=2+ce^6$ by $\sqrt q$ proves the lower bound.
\end{proof}

\begin{thebibliography}{9}
\bibitem{Milman}E. Milman,
\emph{Gaussian Correlation via Inverse Brascamp--Lieb},
Probability Theory and Related Fields (2025), Section 3.3,
\url{https://doi.org/10.1007/s00440-025-01445-x}.

\bibitem{BL}H. J. Brascamp and E. H. Lieb,
\emph{On extensions of the Brunn--Minkowski and Pr\'ekopa--Leindler
 theorems, including inequalities for log concave functions, and with
 an application to the diffusion equation},
Journal of Functional Analysis \textbf{22} (1976), 366--389.

\bibitem{Prekopa}A. Pr\'ekopa,
\emph{On logarithmic concave measures and functions},
Acta Scientiarum Mathematicarum \textbf{34} (1973), 335--343.

\bibitem{Bordenave}C. Bordenave,
\emph{A new proof of Friedman's second eigenvalue theorem and its
extension to random lifts}, arXiv:1502.04482.

\bibitem{MWW}B. D. McKay, N. C. Wormald, and B. Wysocka,
\emph{Short cycles in random regular graphs},
Electronic Journal of Combinatorics \textbf{11} (2004), R66.
\bibitem{DMS}A. Dembo, A. Montanari, and S. Sen,
\emph{Extremal cuts of sparse random graphs},
Annals of Probability \textbf{45} (2017), 1190--1217;
arXiv:1503.03923.
\end{thebibliography}
\end{document}
